\documentclass[output=paper]{langscibook}
\author{Authorname\orcid{}\affiliation{}}
\title{Title}
\abstract{Abstract}
\IfFileExists{../localcommands.tex}{
  \addbibresource{../localbibliography.bib}
  \usepackage{tabularx,multicol}
\usepackage[hyphens]{url}
\urlstyle{same}

\usepackage{langsci-optional}
\usepackage{langsci-lgr}
\usepackage{langsci-gb4e}

\usepackage{soul}
% \usepackage{booktabs}
% \usepackage{geometry}
\usepackage{multirow}
% \usepackage{makecell} %safely breaks a line inside a table cell
% \usepackage{tablefootnote} %footnotes in table captions
\usepackage{enumerate}
\usepackage{tikz}
\usepackage{array}
\usepackage[shortlabels]{enumitem}

%for having several figures on one line, turning words in 90°angles
% \usepackage{graphicx}
\usepackage{subcaption}

% \usepackage{caption} %for making the caption line longer in concrete cases

% \usepackage{rotating} %rotate a whole page; for large tables

\usepackage{fontspec}
\newfontfamily\charis{Charis SIL} %for being able to use the symbol ‿ and a better command of some of the boldfaced diacritics in 01
\newfontfamily\japanesefont{Noto Serif CJK JP} % for Japanese script
% \usepackage{tipa}

% \usepackage{needspace} %to keep exs on one page

\usepackage{xcolor}

\usepackage{pgfplots}
\pgfplotsset{compat=1.18}

%%%%%   AVMs for 02	%%%%%
\usepackage{langsci-avm}
\avmsetup{switch = ?}


\usepackage{langsci-branding}

  %\newcommand*{\orcid}{}


\makeatletter
\let\thetitle\@title
\let\theauthor\@author
\makeatother

\newcommand{\togglepaper}[1][0]{
   \bibliography{../localbibliography}
   \papernote{\scriptsize\normalfont
     \theauthor.
     \titleTemp.
     To appear in:
     E. Di Tor \& Herr Rausgeberin (ed.).
     Booktitle in localcommands.tex.
     Berlin: Language Science Press. [preliminary page numbering]
   }
   \pagenumbering{roman}
   \setcounter{chapter}{#1}
   \addtocounter{chapter}{-1}
}

\newbool{bookcompile}
\booltrue{bookcompile}
\newcommand{\bookorchapter}[2]{\ifbool{bookcompile}{#1}{#2}}



\renewcommand{\thempfootnote}{\fnsymbol{mpfootnote}} % use * for footnotes in float env

% \newcolumntype{L}[1]{>{\raggedright\arraybackslash}p{#1}} %left-aligned (non-justified) text in tables
% \newcolumntype{R}[1]{>{\raggedleft\arraybackslash}p{#1}} % right-aligned (non-justified) text in tables

\definecolor{customgreen}{RGB}{27,158,119}
\definecolor{custompink}{RGB}{231,42,138}
\definecolor{customblue}{RGB}{117,112,179}
\definecolor{customlightblue}{RGB}{143,170,220}
\definecolor{customdarkblue}{RGB}{68,114,196}
\definecolor{customyellow}{RGB}{225,192,0}
\definecolor{customorange}{RGB}{233,137,21}
\definecolor{customgray}{RGB}{229,229,229}


% Left-aligned indexes
% % \makeatletter
% % \patchcmd{\theindex}{\twocolumn}{\raggedright\twocolumn}{}{}
% % \makeatother

%to make Japanese scripts in the bibliography smaller
\newcommand{\japbib}[1]{{\japanesefont\small #1}}

%%%%%%%%%%%%%%%%%%%	MACROS for 02	%%%%%%%%%%%%%%%%

%%% Formatting	%%%%%

\newcommand{\textex}[1]{\textit{#1}} 		 %for formatting linguistic examples in-text%
%command-shift X

\newcommand{\lxm}[1]{\textsc{#1}}  		%for formatting names of lexemes
%command-option-shift L

%%%%%%%%% Abbreviations	%%%%%%

\newcommand{\featname}[1]{\textsc{#1}}
\newcommand{\featspec}[2]{\featname{#1}:{#2}}
\newcommand{\feat}[2]{\textsc{#1}:{#2}}  				% \feat{FEATURE}{value}  sc's
\newcommand{\featnameonly}[1]{\textsc{#1}} 			% \featnameonly{featurename}  sc's

\newcommand{\pr}{$^{\prime}$}

\newcommand{\Corr}{\textit{Corr}}
\newcommand{\propm}{\textit{pm}}

\newcommand{\PC}{Π$_{C}$}

\newcommand{\PF}{Π$_{F}$}
\newcommand{\PR}{Π$_{R}$}

\newcommand{\lex}{$\mathcal{L}$}
 
  %% hyphenation points for line breaks
%% Normally, automatic hyphenation in LaTeX is very good
%% If a word is mis-hyphenated, add it to this file
%%
%% add information to TeX file before \begin{document} with:
%% %% hyphenation points for line breaks
%% Normally, automatic hyphenation in LaTeX is very good
%% If a word is mis-hyphenated, add it to this file
%%
%% add information to TeX file before \begin{document} with:
%% %% hyphenation points for line breaks
%% Normally, automatic hyphenation in LaTeX is very good
%% If a word is mis-hyphenated, add it to this file
%%
%% add information to TeX file before \begin{document} with:
%% \include{localhyphenation}
\hyphenation{
    Al-ex-an-der
    Ber-ber
    chan-ges
    cir-cum-stan-ces
    coin-age
    Fraj-zyngier
    ita-liano
    Ja-pa-nese
    Kö-nig
    Krzy-żanowski
    li-te-ra-tur-no-go
    mi-zen-kei
    par-a-digm
    phe-nom-e-non
    ren-tai-shi
    Slo-vene
    Spen-cer
    Stoj-kov
    Ver-bi-ca-re-se
    wide-spread
}

\hyphenation{
    Al-ex-an-der
    Ber-ber
    chan-ges
    cir-cum-stan-ces
    coin-age
    Fraj-zyngier
    ita-liano
    Ja-pa-nese
    Kö-nig
    Krzy-żanowski
    li-te-ra-tur-no-go
    mi-zen-kei
    par-a-digm
    phe-nom-e-non
    ren-tai-shi
    Slo-vene
    Spen-cer
    Stoj-kov
    Ver-bi-ca-re-se
    wide-spread
}

\hyphenation{
    Al-ex-an-der
    Ber-ber
    chan-ges
    cir-cum-stan-ces
    coin-age
    Fraj-zyngier
    ita-liano
    Ja-pa-nese
    Kö-nig
    Krzy-żanowski
    li-te-ra-tur-no-go
    mi-zen-kei
    par-a-digm
    phe-nom-e-non
    ren-tai-shi
    Slo-vene
    Spen-cer
    Stoj-kov
    Ver-bi-ca-re-se
    wide-spread
}
 
  \togglepaper[1]%%chapternumber
}{}

\begin{document}
\maketitle 
%\shorttitlerunninghead{}%%use this for an abridged title in the page headers


\textbf{Diachronic} \textbf{paths} \textbf{to} \textbf{uninflectedness} \textbf{in} \textbf{South} \textbf{Slavic}\footnote{The research reported here was funded by the Leverhulme Trust, grant number RPG-2020-078, as part of project “Declining case: inflectional loss in progress” (https://www.smg.surrey.ac.uk/projects/declining-case-inflectional-loss-in-progress). The authors are grateful to two anonymous reviewers whose comments and suggestions led to significant improvements in the paper.}

Alexander Krasovitsky, Maria Kyuseva, Matthew Baerman, Greville G. Corbett

\textit{We argue that the ongoing loss of case in Serbian and Bulgarian dialects has led to a split system with competing co-grammars: one where nouns inflect and another where they do not. This division between "forerunners" (without inflection) and "underachievers" (with inflection) results in partial rules of different kinds. There are three types of splits depending on the dialect. First, a split between morphological classes, where some retain case distinctions while others do not. Second, classes that retain inflection may still exhibit uninflectedness in specific constructions, such as partitive constructions. Third, an otherwise inflecting morphological class may split on a semantic basis, such as human vs. non-human. These asymmetries in morphological change lead to the rise of partial synchronic rules, whether morphological, syntactic, or semantic. These rules persist over time and evolve as historical processes move into new phases. By studying these rules in contemporary dialects and arranging them chronologically, we can uncover the detailed stages of this historical process.}

\begin{itemize}
\item \begin{styleListParagraph}
\textbf{Introduction}
\end{styleListParagraph}
\end{itemize}

The loss of inflectional case marking is an ongoing process in the dialect continuum that stretches across Serbian and Bulgarian linguistic territory. This has led to substantial variation in the dialect case systems of the region. We observe the first signs of this historical change in systems with rich case distinctions, where case marking is obligatory on all classes of nominals, as in the western part of the continuum. Further east, we find dramatically reduced case paradigms, with just a few distinct case forms, attested with only part of the lexicon, while the rest of the lexicon has lost the ability to inflect for case. At the same time, in these dialects we find competing rules of case marking, which appear to reflect different historical periods in the breakdown of the case system. The result is either competition between more or less distinct inflection (towards the west, in Serbian dialects), or competition between inflection and uninflectedness (towards the east, in Bulgarian dialects). We investigate factors which underlie this competition and show how they vary at different stages of this case decline. This provides us with an insight into the historical mechanisms underlying case loss, and to the conditions that may foster this change.

While our analysis has confirmed the relevance of the factors previously proposed to explain variable case marking synchronically and diachronically (e. g. \citealt{Moravcsik1978}; \citealt{Bossong1991}; \citealt{König2008}; \citealt{DalrympleNikolaeva2011}; \citealt{Cristofaro2013}; Kurumada, \& \citealt{Jaeger2015}; \citealt{Witzlack-MakarevichSeržant2018}; \citealt{SeržantWitzlack-Makarevich2018}), the data obtained from Serbian and Bulgarian dialects shed new light on these factors and show that they do not in fact operate in the way that has previously been suggested. 

First, extensive typological data show \textstylecfi{explicit case marking} is used in environments where it is essential to avoid ambiguity (e.g. between subjects and objects) and is omitted where syntactic roles are transparent \citep{Comrie1977}. For example, it has been shown that the omission of morphological case marking goes hand-in-hand with the use of adpositions, which to a certain extent duplicate the information conveyed by case inflection. This has been shown to play a role in diachronic processes affecting case systems: the information redundancy of case inflection in adpositional constructions may make these constructions more prone to the loss of morphological case than those without adpositions, where case inflection provides critical information about syntactic roles (\citealt{HewsonBubenik2006}). In the dialects of East Serbia, while the use of adpositions (prepositions) correlates with a decline in case distinctions, the dialects display a pattern which is the opposite of what is predicted by the above principle: the original case marking is better preserved precisely when governed by a preposition, and more prone to loss elsewhere.

Second, pragmatic requirements, in particular discourse salience, have been proved to be another important factor in variable case marking cross-linguistically. To mention just two examp les, \citet{McGregor2018} and \citet{Montaut2018} present somewhat similar situations in unrelated languages (Southern African Khoisan and Dravidian respectively) where the rise of object case marking led to pragmatically motivated choices between overt case marking (accusative) and zero marking: prominent objects receive overt case marking while backgrounded objects do not. In Bulgarian border-region dialects, where just two cases have been preserved (nominative and accusative), explicit case marking on non-subject nouns (accusative) is also conditioned by pragmatic factors. However, it is not associated with pragmatic salience, but rather the opposite:  explicit (accusative) marking on non-subjects is more frequent in pragmatically less prominent parts of an utterance (e.g. containing background information). Pragmatic prominence instead tends to trigger nominative marking on non-subject nouns, which makes them morphologically indistinguishable from subjects. 

We present two case studies that show how different factors shape the transition from a larger to a smaller case system, and further to uninflectedness. After the essential background information in \sectref{sec:key:2}, we consider the role of structural factors, which are dominant in Serbian dialects (\sectref{sec:key:3}), and analyse the impact of   pragmatic conditions, which are dominant in Bulgarian dialects under study (\sectref{sec:key:4}). We suggest that this shift in prominence is due to the diachronic stage each variety has reached: structural factors are prominent at earlier stages of case decline and pragmatic factors at an advanced stage.

\begin{itemize}
\item \begin{styleListParagraph}
\textbf{Background}
\end{styleListParagraph}
\end{itemize}

South Slavic varieties inherited from Proto-Slavic a case system with six morphological cases (plus the vocative whose status as case is debatable\footnote{Vocative forms continue to be widely used in contemporary Serbian and Bulgarian, but do not express grammatical relationships \citep[8]{Blake2004} and do not interact with the types of configurations we have investigated.} ) and six major inflection classes (\citealt{Mirčev1958}: 147-160, \citealt{Slavova2017}: 127-178, \citealt{Popović1955}: 115). The further history of case largely correlates with the division of the South Slavic territory into eastern and western zones. In the western zone (contemporary Serbian, Croatian, Bosnian, Montenegrin, Slovenian) the original Proto-Slavic case system has been largely preserved. In the eastern zone (contemporary Bulgarian and Macedonian), most varieties have seen the reduction and loss of nominal case marking and the rise of analytical constructions (e.g. prepositional constructions), as can be seen by contrasting the data from standard Bulgarian and standard Serbian in \tabref{tab:key:1}\footnote{Bulgarian and Macedonian demonstrate an asymmetry between nominal and pronominal systems with respect to case marking. Nouns and adjectives do not inflect for case (in standard varieties and in a significant number of dialects). Personal pronouns distinguish full and clitic forms (the choice being syntactically determined). Full personal pronouns retain a two-case distinction, nominative vs. accusative, the latter being used in all non-subject syntactic roles: \textbf{\textit{tja}} (3SG.F.NOM) \textit{obi}\textrm{\textit{č}}\textit{a} (’she loves’) vs. \textit{obi}\textrm{\textit{č}}\textit{am} \textbf{\textit{neja}} (3SG.F.ACC) (‘I love her’) vs. \textit{kazax na} \textbf{\textit{neja}} (3SG.F.ACC) (‘I told her’, literally ‘to her’). Clitic personal pronouns are used in non-subject roles and distinguish between accusative and dative: \textit{obi}\textrm{\textit{č}}\textit{am} \textbf{\textit{ja}} (3SG.F.ACC) vs. \textit{kazax} \textrm{\textbf{\textit{ì}}} (3SG.F.DAT)  (‘I told her’).}.


\begin{tabularx}{\textwidth}{XXX}

\lsptoprule

Standard Serbian: inflected forms & Standard Bulgarian: prepositional phrases with an uninflected form & Translation\\
ovo je \textbf{Kipar} (nominative)

idu na \textbf{Kipar} (accusative)

pomažu \textbf{Kipru} (dative)

stanovništvo \textbf{Kipra} (genitive)

žive na \textbf{Kipru}  (locative)

upravljaju \textbf{Kiprom} (instrumental) & tova e \textbf{Kipâr}

otivat \textbf{na} \textbf{Kipâr}

pomagat \textbf{na} \textbf{Kipâr}

 naselenieto \textbf{na} \textbf{Kipâr}

živejat \textbf{na} \textbf{Kipâr}

upravljavat  \textbf{Kipâr} & ‘This is Cyprus’ 

{}'They go to Cyprus'

{}'They help Cyprus'

{}'The population of Cyprus'

‘They live in Cyprus’

‘They govern Cyprus’\\
\lspbottomrule
\end{tabularx}
%%please move \begin{table} just above \begin{tabular
\begin{table}
\caption{Nominal inflection in Serbian and corresponding forms in Bulgarian}
\label{tab:key:1}
\end{table}

In between these two types there is a substantial transitional zone. A characteristic feature of these dialects is the competition between these two strategies, as seen both in Serbian \REF{ex:key:1} and Bulgarian dialects \REF{ex:key:2}\footnote{For interlinear glosses, we follow the conventions of the Leipzig Glossing Rules (\url{https://www.eva.mpg.de/lingua/pdf/Glossing-Rules.pdf}) with the following additions: AOR – aorist, EXIST – existential, IMPF – imperfect.} .

\textit{1.  Zove}  \textit{ga}  \textit{otac}  \textbf{\textit{telefon-om}}  \textit{/}  \textbf{\textit{na}}  \textbf{\textit{telefon}}

  call[\textsc{prs.3sg]}  \textsc{3sg.m}.\textsc{acc}  father[\textsc{nom]}  phone-\textsc{ins.sg}  on  phone[\textsc{acc](=nom)}  ‘The father calls him by phone’    (\citealt{Miloradović2002}: 183; 202)  

2.  \textbf{\textit{Petr-a} }   / \textbf{\textit{Petâr} }   \textit{go}    \textit{njama}

Peter-\textsc{sg.acc} /  Peter[\textsc{sg.nom]}  \textsc{3sg.m}.\textsc{acc}  \textsc{exist.neg}\\
‘Peter is absent’            \citep[163]{Stojkov1981}

The dialects in this transitional area correspond to different postulated chronological stages of the decline of the Proto-Slavic case system. A group of dialects in the west of the transitional area preserve all six original Proto-Slavic cases (plus vocative), but speakers regularly replace three of them with the accusative form, which serves as a general oblique case (\tabref{tab:key:2}, Transitional system 1). Thus, the original six-case system competes in these dialects with an innovative three-case system, consisting of nominative, dative, and accusative only (\citealt{Miloradović2003}; \citealt{Simić1972}). The further east one goes, the greater are the reductions. The dative case is the next to be replaced (Transitional system 2). Here, the more conservative three-case system competes with a two-case system, consisting of nominative and accusative (\citealt{Belić1905}, \citealt{Ivić1956}). Finally, in the east of the transitional area, the accusative itself starts to yield to a single uninflected form (Transitional System 3). Here, the two-case system coexists with a system without case (\citealt{Stojkov1975,Stojkov1981}). Thus, three structural dialect types can be distinguished within the transitional area, according to the size of the case paradigm: six-, three- and two-case systems, as summarised in \tabref{tab:key:2}.  

  
%%please move the includegraphics inside the {figure} environment
%%\includegraphics[width=\textwidth]{figures/a5Krasovitskyetalrevisedversion-img001.png}
 

%%please move \begin{table} just above \begin{tabular
\begin{table}
\caption{Case systems in the Serbian-Bulgarian dialect continuum\footnote{In distinguishing between the largely syncretic dative and locative cases in Serbian, we follow the Slavic tradition of case description (Piper, \citealt{Klajn2013}, \citealt{Miloradović2003}). The non-syncretic locative endings were used at least until the end of the 15th century \citep{Belić1999}, while the loss of case started sometime before the 12th century \citep{Belić1905}. It thus that when the process of the case loss started, the dative and locative were still morphologically distinct cases. An additional argument in favour of differentiating between the two cases comes from the fact that it facilitates the description of the process of the case loss. The original locative case is frequently replaced by the general oblique in the whole transitional area, including the western region where it is the only case to be replaced. The original dative case, on the contrary, is preserved in most Serbian dialects with the exception of the eastern zone where the only remaining cases are nominative and accusative \citep{Sobolev1991}.}}
\label{tab:key:2}
\end{table}

The chronology that this transition reflects has been reconstructed primarily on the basis of liturgical texts produced within various regional written traditions in the South Slavic territory from the 10th century onwards (\citealt{Tsonev1984}; \citealt{Duridanov1956,Duridanov1958}; \citealt{Rusek1964}; \citealt{Češko1970}; \citealt{Belić1905}) and, for later periods, on the basis of clerical and household documents \citep{Bernštejn1948}. 

It is believed that the decay of the case system originated in the South Slavic eastern zone (Bulgarian and Macedonian) around the 10th century \citep[25]{Duridanov1958}. By the end of the 14th century morphological case marking on nominals was completely lost in at least in some of the dialects spoken in the eastern zone (Bulgarian border-region dialects) \citep[11]{Češko1970}. Even earlier, by the end of the 13th century, case loss was at an advanced stage in a number of Macedonian dialects \citep{Rusek1964}. In Serbian, the process of case decline started sometime before the 12th century and affected the South-East dialects \citep{Belić1905}. It did not, however, lead to the complete loss of case distinctions: rather it resulted in different types of reduced systems.

Historical evidence points to the following processes involved in this change: the spread of prepositional constructions replacing bare noun phrases; the redistribution of functions between indirect cases and displacement of less productive indirect cases; the spread of accusative forms in the function of an oblique case generalising across original indirect cases\footnote{The chronology of the displacement of individual indirect cases varies. The most stable, as dialectological data of the Serbian language show (\citealt{Belić1905}, \citealt{Miloradović2003}) is the dative case.  The historical data from Bulgarian also testify to this: the dative forms of nouns were found in Bulgarian texts of the 17th, 18th and even 19th centuries, while the genitive, instrumental and local cases had already fallen out of use in Bulgarian dialects in the 13th-14th centuries (\citealt{Mirčev1958}: 255-260; Če\textrm{š}ko 1970: 301-310).}; and the spread of uninflected forms replacing this oblique case.

\begin{itemize}
\item \begin{styleListParagraph}
\textbf{Constructional} \textbf{variation:} \textbf{the} \textbf{loss} \textbf{of} \textbf{the} \textbf{genitive} \textbf{in} \textbf{Serbian} \textbf{dialects}
\end{styleListParagraph}
\end{itemize}

The loss of case in European languages has been reported to go hand in hand with the increased use of prepositions (\citealt{Blake2004}; \citealt{HewsonBubenik2006}; \citealt{Kulikov2006}). \citet[178]{Blake2004} illustrates this with an Old English version of Luke 15:15 \REF{ex:key:3a} and Wyclif’s fourteenth-century translation \REF{ex:key:3b} (glossing and translation are as in \citealt{Blake2004}):

3a. \textit{He    folgode       an-um        burg-sitt-end-um      menn                  thaes          rices}

he     followed  one-\textsc{dat}  town-dwell-ing-\textsc{dat}  man.\textsc{dat}            that.\textsc{gen}     land.\textsc{gen}

3b. \textit{He clevede} [=attached himself] \textit{to oon of the citizens of that contré}

‘He hired himself out to one of the citizens of that country.’

In \REF{ex:key:3a}, the directional (‘to the citizen’) and the possessive (‘of this country’) relationships are expressed via inflection, while in \REF{ex:key:3b}, they are expressed by prepositional constructions. The relationship between the two processes, however, is not clear. Did the spread of prepositional constructions cause case to disappear as it became semantically redundant (as claimed in \citealt{Hjelmslev1935}; \citealt{HewsonBubenik2006})? Or was it, by contrast, a repair strategy necessary after the loss of explicit case marking, as proposed, for example, in \citet{Skautrup1944}? 

The variation found in Serbian dialects allows us to investigate the interplay of inflectional case marking and prepositional constructions using synchronic data. We focus here on the genitive, which is in the process of being replaced by the accusative. Specifically, we ask: do prepositional constructions particularly favour this replacement? This is what an informativeness principle would predict: since the preposition provides semantic context, genitive case marking is redundant and the more general accusative case can be used without any loss of information (cf. \citealt{Comrie1977,Comrie1989}). Our data, however, show the opposite tendency: prepositional constructions favour the \textit{retention} of the genitive, while its replacement by the accusative is more advanced in non-prepositional contexts.  That means that in Serbian dialects, some principle other than informativeness is at play. At this point in our investigations we are unable to identify what this could be, so our goal here is to provide evidence which might spur further research. 

In \sectref{sec:key:3.1}, we introduce the data used in this case study. \sectref{sec:key:3.2} outlines the use of the genitive forms in Serbian and discusses the results of previous research on the variation between genitive and accusative in South-East Serbian dialects. Sections 3.3 and 3.4 present the patterns of variation found in our data; \sectref{sec:key:3.5} discusses the theoretical implications of our findings.

\textbf{3.1.} \textbf{Data}

The data for this study come from sociolinguistic interviews with speakers of Serbian dialects located in or near Kosovo. These come from two sources: the archive of the Institute for Balkan Studies and our own fieldwork. The archive provided a subset of recordings collected in the project “Research of Slavic Vernaculars in Kosovo and Metohija” (2002 – 2003) led by the Institute for the Serbian Language of SASA and financed by UNESCO. Our data come from fieldwork in the municipality of Brus, Central Serbia, conducted by Maria Kyuseva within the project “Declining case: inflectional loss in progress” in 2022. 

These data were collected in a uniform fashion that conforms to the Serbian dialectological tradition. The participants were native speakers of local dialects, aged between 70 and 81 years, who had been long-term residents of the region and had engaged in traditional activities throughout their lives, such as cattle breeding or agriculture. This sampling ensured the minimal interference of the high-prestige standardised variety of Serbian. The topics of the interviews included history, tradition, culture, crafts, cuisine, everyday life, and biographical stories.

The total corpus compiled from these two sources contains over 274,800 tokens, 30 hours and 36 minutes of recordings. Overall, twenty-three speakers from eleven villages were recorded. \figref{fig:key:1} shows the locations of the villages, which are colour-coded according to dialect. The Zeta-South Sandžak dialect is marked in green, the Kosovo-Resava dialect is marked in pink, and the Prizren-South Morava dialect is marked in purple.

   
%%please move the includegraphics inside the {figure} environment
%%\includegraphics[width=\textwidth]{figures/a5Krasovitskyetalrevisedversion-img002.png}
 

\begin{figure}
\caption{Locations of the interviews in Serbia}
\label{fig:key:1}
\end{figure}

%%[Warning: Draw object ignored]
%%[Warning: Draw object ignored]
%%[Warning: Draw object ignored]
      Zeta-South Sandžak  Kosovo{}-Resava  Prizren-Timok

With respect to case marking, these three dialects have the same system (Transitional System 1 in \tabref{tab:key:2}): all six original cases are preserved to various degrees, but three of them (genitive, locative, and instrumental) can be replaced by accusative forms. The dialects differ in the frequency of such replacements. Zeta-South Sandžak is the most conservative: it has a stable six-case system, and the three peripheral cases are only occasionally replaced by the accusative. Kosovo-Resava represents a more advanced stage of case loss: the genitive, locative, and instrumental cases frequently give way to the accusative. Finally, Prizren-South Morava shows the most advanced stage: the locative and instrumental are virtually non-existent (occurring with only a handful of lexemes), and the use of genitive is highly restricted. We compare the usage of genitive forms in these three dialects and project the synchronic tendencies onto the diachronic plane. This allows us to draw conclusions regarding the historical dynamics involved in the decline of this case.

\textbf{3.2.} \textbf{Genitive} \textbf{case} \textbf{in} \textbf{Serbian}

The Serbian genitive is highly polysemous. Its typical function is to mark a relationship between two entities, such as possessor-possessee, part-whole, and source-goal. Additionally, it can denote time, reason, goal, and manner. From the point of view of syntax, it has a wide distribution. It can be used in adnominal, adjectival, and verbal constructions, as well as in prepositional constructions and constructions with a quantifier. Out of all the cases, the genitive occurs with the widest inventory of prepositions \citep{Feleško1955}.

Various studies of individual dialects of South-East Serbia have already established the heterogeneous behaviour of the genitive in different constructions. For example, \citet{Miloradović2003} analyses the competition between the genitive and accusative in the Kosovo-Resava dialect of the Paraćin region in Central Serbia. She observes that while the genitive is dominant in prepositional constructions, its use is highly limited elsewhere. In verbal constructions, the genitive is replaced by the accusative \REF{ex:key:4a}, in adnominal constructions – by a prepositional phrase \REF{ex:key:4b}, and in adjectival constructions, it can be replaced either by the accusative or by a prepositional phrase \REF{ex:key:4c}.

4a.   \textit{nije        bi-l-o} \textbf{\textit{vod-e}} {}-\textit{>  nije           bi-l-o} \textbf{\textit{vod-u}}

 \textsc{aux.neg.3sg}   be-\textsc{ptcp-n.sg}   water-\textsc{gen.sg}             \textsc{aux.neg.3sg}  be-\textsc{ptcp-n.sg}  water-\textsc{acc.sg}

         ‘There was no water’

4b.    \textit{kraj} \textbf{\textit{mačug-e}} {}->         \textit{kraj} \textbf{\textit{od}} \textbf{  \textbf{mačug-e}}

         end\textsc{[nom.sg]}       cane-\textsc{gen.sg}  end[\textsc{nom.sg}]        of        cane-\textsc{gen.sg}

         ‘Cane end’

4c.    \textit{pun} \textbf{\textit{vod-e} } {}->         \textit{pun}                  \textbf{\textit{vod-u}}

   full[\textsc{m.nom.sg}]  water-\textsc{gen.sg}  full[\textsc{m.nom.sg]}    water-\textsc{acc.sg}

      \textit{pun} \textbf{\textit{s}} \textbf{\textit{vod-u}}

 full[\textsc{m.nom.sg]}  with  water-\textsc{acc.sg}

‘[It is] full of water’

\citet{Stanković2022} analyses the speech of the younger population in the town of Vranje in South Serbia, where the Prizren-South Morava dialect is spoken. This dialect has a three-case system consisting of the nominative, accusative, and dative. The speech of the younger generation, however, is influenced by the standard dialect. Therefore, it is susceptible to the penetration of other cases found in the standard, especially the genitive. Stanković observes that the possibility of using the genitive depends on the construction. Focusing on the non-prepositional uses, she finds that adnominal constructions and constructions with a numeral permit the occasional appearance of the genitive case (5a-b), while constructions where the complement depends on a verb or a predicative adjective do not (5c-d):

5a.   \textit{udic-a               za} \textbf{\textit{pecanj-e}                \textbf{rib-a}}

 rod-\textsc{nom.sg}     for     catching-\textsc{acc.sg}      fish-\textsc{gen.pl}

       ‘Fishing rod’

5b.   \textit{ima-m}   \textbf{\textit{pet}    }  \textbf{\textit{različit-ih}                        }  \textbf{\textit{igar-a}}

       have-\textsc{prs.1sg}           five     different-\textsc{gen.pl}           game-\textsc{gen.pl}

         ‘I have five different games’

5c.     \textit{nije} \textbf{\textit{je-o}                            \textbf{čokolad-u}  } \textit{/    *čokolad-e}

       \textsc{aux.neg.3sg}    eat-\textsc{ptcp.m.sg}           chocolate-\textsc{acc.sg}  / *chocolate-\textsc{gen.sg}

       ‘He didn’t eat the chocolate’

5d.   \textit{dvorišt-e} \textbf{\textit{pun-o}}  \textbf{\textit{s}}  \textbf{\textit{igračk-e} }  \textit{/  *pun-o  igračak-a}  yard-\textsc{nom.sg}  full-\textsc{n.nom.sg}\textbf{  }with    toy-\textsc{acc.pl}    / *full-\textsc{n.nom.sg}  toy-\textsc{gen.pl}

       ‘The yard is full of toys’

In our study, we bring these individual observations together by comparing the use of the genitive case forms in different constructions across the three South-East Serbian dialects: Zeta-South Sandžak, Kosovo-Resava, and Prizren-South Morava.

\textbf{3.3.}  \textbf{Constructions} \textbf{with} \textbf{the} \textbf{genitive} \textbf{in} \textbf{South-East} \textbf{Serbian}

We focus on three constructions with the genitive case: prepositional constructions, constructions with a quantifier, and adnominal constructions. In South-East Serbian dialects, these constructions allow both the original genitive and the innovative accusative forms:

6. Prepositional construction:

6a.    \textit{napravljen-i}   \textbf{\textit{od}}   \textbf{\textit{drv-a}}

     made-\textsc{m.nom.pl}          from                       wood-\textsc{gen.sg}

 ‘[They are] made from wood’

6b.   \textit{ima                        i} \textbf{\textit{od}} \textbf{\textit{drv-o}}

\textsc{exist.prs}           and     from   wood-\textsc{acc.sg}

‘There is [some of it] from wood as well’ 

7. Construction with a quantifier:

7a.   \textit{malo} \textbf{\textit{trav-e}} \textit{gi          stavi-m}

          a.little      grass\textsc{{}-gen.sg}   3\textsc{pl.dat}  put-\textsc{prs.1sg}

‘I put a bit of grass for them’

7b. \textit{i          ondak         turi                 malo} \textbf{\textit{trav-u}}

       and  then            put\textsc{[imp}]        a.little              grass-\textsc{acc.sg}

    ‘And then put a little bit of grass’  

8. Adnominal construction:

8a.        \textit{tamo         odnese-š                         čaš-u} \textbf{\textit{vin-a}}

       there        take.away-\textsc{prs.2sg}       cup-\textsc{acc.sg}       wine-\textsc{gen.sg}

       ’You take there a cup of wine\textsc{’}

8b.         \textit{i    uzme-š              čaš-u} \textbf{\textit{vin-o}}

    and  take-\textsc{prs.2sg}    cup-\textsc{acc.sg}       wine\textbf{{}-}\textbf{\textsc{acc.sg}}

   ‘And you take a cup of wine’

The overall number of examples in the dataset is 960. The prepositional construction is represented by 676 examples, the adnominal construction – by 169 examples, and the construction with a quantifier – by 115 examples. The ratio represents the occurrence of these constructions in spontaneous speech. For the prepositional construction, we include only phrases with the preposition \textit{od} ‘from, of’, which is the most frequent preposition that governs the genitive. Constructions with a quantifier include the following quantifiers: \textit{mnogo} ‘a lot’, \textit{malo} ‘a little’, \textit{lek} ‘very little’, \textit{pomalo} ‘a little bit’, \textit{puno} ‘much’, \textit{više} ‘more’, and \textit{manje} ‘less’. Finally, adnominal constructions include a variety of governing nouns, the most frequent meanings in the sample are: “quantity” (a cup of tea, a bowl of rice), and “part-whole” (top of the table, edge of the village). True possession (a son of my brother), otherwise typical for the construction, is rare in our examples. It is usually expressed in the dialects by alternative means of expression, such as possessive adjectives and constructions with the dative case.

The bar chart (\figref{fig:key:2}) shows the variation between the original genitive and the innovative accusative in the three dialects.

  
%%please move the includegraphics inside the {figure} environment
%%\includegraphics[width=\textwidth]{figures/a5Krasovitskyetalrevisedversion-img003.png}
 

\begin{figure}
\caption{Competing strategies of case marking in South-East Serbian dialects}
\label{fig:key:2}
\end{figure}

As the figure illustrates, in Zeta-South Sandžak (the leftmost section of the bar chart), the most conservative dialect in the sample, the genitive prevails in all three constructions. It appears in 90\% or more of the examples in the prepositional and adnominal constructions and is used somewhat less in the construction with a quantifier. As the ANOVA\footnote{The ANOVA analysis of the corpus data here and below was carried out by Dr Alexander Stewart (University of St Andrews, UK).} test shows, there is no significant effect of the construction on the case in this dialect, with p=0.2429 and F = 1.4284. The difference between the three constructions is more prominent in the Kosovo-Resava dialect, which has a more advanced stage of case decline. Here, the genitive still dominates with 90\% of occurrences in the prepositional construction, it concedes slightly in the adnominal construction, and its use falls dramatically in the construction with a quantifier, where it is no longer a dominant strategy. The ANOVA test shows a significant effect of the construction on the case in this dialect with p=3.3092e-18 and F=44.1388. Finally, the most innovative Prizren-South Morava dialect (the rightmost section of the bar chart) shows yet another picture. Here, the use of accusative is the dominant strategy in both the adnominal construction (65\% of examples) and the construction with a quantifier (92\% of examples), and the only construction where the genitive appears in a slight majority of examples (56\%) is the prepositional construction. Again, ANOVA shows a significant effect of the construction with p=1.1018e-06 and F=14.2548. (The null hypothesis for each of three dialects assumed no effect of the construction on case selection.)

These data show that the use of the genitive versus the accusative in the Serbian dialects under study is at least partly conditioned by the type of construction. The exact conditions depend on the dialect, The exact conditions depend on the dialect, or if we accept the diachronic interpretation of the data, on the stage of case decline. The general tendency is the following:  constructions with a quantifier are the most susceptible to the penetration of the innovative accusative, adnominal constructions come second, and prepositional constructions preserve the original genitive the longest. This effect is clear when we look at the distribution of cases in individual dialects, but it can also be seen if we trace the linguistic behaviour of the constructions across dialects. The use of the genitive case consistently declines in each construction as we move across the transitional zone from west to east. For each construction, we find a significant effect of the dialect: with p=2.1082e-25 and F=61.8976 (prepositional construction), p=9.1127e-10 and F=23.6043 (adnominal construction), and p=1.1889e-06 and F=15.4478 (construction with a quantifier).

\textbf{3.4.} \textbf{Morphologically} \textbf{conditioned} \textbf{variation} \textbf{in} \textbf{Serbian} \textbf{dialects}

Another factor that conditions variation between the original genitive and innovative accusative case is the inflection class.  This plays a role in the whole dialectal continuum, but in our data it is especially pronounced in the Prizren-South Morava dialect. 

In standard Serbian, nouns cluster into four inflection classes (\tabref{tab:key:3}). The paradigms exhibit a high degree of syncretism (see the comment in footnote 5 concerning dative – locative syncretic forms).


\begin{tabularx}{\textwidth}{XXXXXXXXX}
 & \multicolumn{2}{c}{{ Inflection class I:}
\lsptoprule

{ prostor ‘space’}

 masculine} & \multicolumn{2}{c}{{ Inflection class II:}

{ kuća ‘house’}

 feminine (some masculine)} & \multicolumn{2}{c}{{ Inflection class III:}

{ stvar ‘thing’}

 feminine} & \multicolumn{2}{c}{{ Inflection class IV:}

{ polje ‘field’}

 neuter}\\
& SG & PL & SG & PL & SG & PL & SG & PL\\
NOM & prostor & prostor-i & kuć-a & kuć-e & stvar & stvar-i & polj-e & polj-a\\
ACC & prostor & prostor-e & kuć-u & kuć-e & stvar & stvar-i & polj-e & polj-a\\
GEN & prostor-a & prostor-ā & kuć-e & kuć-ā & stvar-i & stvar-ī & polj-a & polj-ā\\
DAT & prostor-u & prostor-ima & kuć-i & kuć-ama & stvar-i & stvar-ima & polj-u & polj-ima\\
LOC & prostor-u & prostor-ima & kuć-i & kuć-ama & stvar-i & stvar-ima & polj-u & polj-ima\\
INS & prostor-om & prostor-ima & kuć-om & kuć-ama & stvar-ju & stvar-ima & polj-em & polj-ima\\
\lspbottomrule
\end{tabularx}
%%please move \begin{table} just above \begin{tabular
\begin{table}
\caption{Inflection classes in Serbian.}
\label{tab:key:3}
\end{table}

In South-East Serbian dialects, there is a tendency for nouns of inflection class III to adopt the forms (and the gender) of inflection class I. For this reason, we do not have enough examples to be able to draw reliable generalizations about this class, and we exclude it from further analysis.  The bar chart below shows the variation between the genitive and the accusative for the remaining three inflection classes in the Prizren-South Morava dialect:

  
%%please move the includegraphics inside the {figure} environment
%%\includegraphics[width=\textwidth]{figures/a5Krasovitskyetalrevisedversion-img004.png}
 

\begin{figure}
\caption{Prizren-South Morava dialect: original genitive vs. innovative accusative in different inflection classes}
\label{fig:key:4}
\end{figure}

The choice of case is clearly conditioned by inflection class: class II favours the genitive form, while class I and class IV favour the accusative form. This pattern can be illustrated with examples of the same meaning of the genitive (material) being expressed differently depending on the inflection class:

9a.   se         ispred-e            \textbf{od}                   \textbf{vun-e}

\textbf{  }  \textsc{refl}      spin-\textsc{prs.3sg}   from           wool(II)-\textsc{gen.sg}

       ‘It is spun from wool’

9b. pravi-l-i             \textbf{od}          \textbf{drv-o}

      make-\textsc{ptcp-m.pl}  from    wood(IV)\textsc{{}-acc.sg}

      ‘They used to make it from wood’

In \REF{ex:key:9a}, the noun belongs to class II (\textit{vuna} ‘wool’) and takes the genitive; in \REF{ex:key:9b} the noun belongs to class IV (\textit{drvo} ‘wood’) and takes the accusative.  This tendency to preserve the original case form in class II is even more pronounced in Bulgarian border-region dialects, where case marking is only possible with class II nouns (see \sectref{sec:key:4}).

\textbf{3.5.} \textbf{Discussion:} \textbf{structural} \textbf{factors} \textbf{in} \textbf{case} \textbf{decline}

These data illustrate both syntactic and morphological conditions on variation between the genitive and accusative. In terms of syntax, certain constructions favour the use of the accusative over the genitive, in particular prepositional phrases. In terms of morphology, certain inflection classes are more likely to adopt the innovative accusative, in particular class II. 

The tendency of prepositional phrases to preserve the original case marking contradicts commonly held assumptions about the role played by prepositions in case loss (\citealt{Blake2004}; \citealt{HewsonBubenik2006}). Since genitive case marking in these prepositional constructions is semantically redundant, it could be omitted without incurring any ambiguity. One might assume then that prepositional constructions would be the first to lose the genitive. The fact that our data show the opposite means that the preservation of the genitive is not driven by the need to disambiguate. We propose here a possible alternative motivation.

Following \citet{Norde2002}, we distinguish two types of less-effort strategies for grammatical change: a speaker-oriented and a hearer-oriented strategy. These can be thought of as two opposing forces that shape any language change. The speaker “wants” to make less effort in encoding the message, and the hearer “wants” to make less effort in decoding the message. The informativeness principle aligns with a hearer-oriented strategy. In this instance, the genitive should first be replaced in prepositional constructions, as it is not needed for the hearer to be able to correctly interpret the message, while it has a function in other constructions. \textcyrillic{А} speaker-oriented strategy would be one that favours ease of production. The genitive case after the preposition \textit{od} ‘from’ is syntactically determined: \textit{od} governs the genitive, creating a high degree of predictability: every time the preposition \textit{od} ‘from’ is used, it can be expected to be followed by the genitive. In adnominal constructions and constructions with a quantifier, there is no such syntactic determination, in the sense that the nouns, such as ‘cup’ or ‘edge’, do not obligatorily govern a noun phrase, and quantifiers can similarly be used on their own. In this way, the syntactic information supports the case morphology and contributes to the preservation of the original case forms. Crucially, there is no gain for the hearer in preserving the case forms in prepositional constructions while omitting them in non-prepositional constructions. Rather, it is in the speaker’s interest to keep an inflected form in a context that always requires it.\textstyleFootnoteSymbol{} 

With these results in mind, we turn to the contrasting patterns of Bulgarian border-region dialects.

\textbf{4.} \textbf{Competing} \textbf{case} \textbf{marking} \textbf{in} \textbf{Bulgarian} \textbf{border-region} \textbf{dialects:} \textbf{pragmatic} \textbf{factors}

In Bulgarian border-region dialects, case marking is lost in all inflection classes of nouns except class II. In this class, singular nouns retain two case forms, nominative and accusative; plural forms are not inflected for case. The nominative marks nouns in the subject position and the accusative marks nouns in non-subject positions. Nouns in classes I, III and IV have retained only one (historically nominative) morphological form for each number value (singular and plural) which is used both in the subject and any non-subject position. This is illustrated in examples \REF{ex:key:10}, \REF{ex:key:11}, \REF{ex:key:12} and \REF{ex:key:13} from our corpus (data represent various locations in the municipalities of Belogradchik and Trân). 

Inflection class I, masculine. Caseless (historically nominative)


\begin{tabularx}{\textwidth}{XXXX}

\lsptoprule
10a. & \textit{Dojde}     \textbf{\textit{doktor.}}

Come[\textsc{pst.3sg]}  doctor[\textsc{sg]}

‘A doctor came.’ & subject NP & Filipovtsi,                                                                                                            \\
Trân\\
10b. & \textit{Izvika-me} \textbf{\textit{doktor.}}

call-\textsc{prs.1pl}  doctor[\textsc{sg]}

‘We call a doctor.’ & non-subject NP & Filipovtsi,                                                                                                              \\
Trân\\
10c. & \textit{Tija                    babičk-i}    \textit{kak    da       otiva-t} \textbf{   }

this [\textsc{nom.pl.f}]   elderly.woman-\textsc{nom.pl}     how   \textsc{comp}  go-\textsc{prs.3pl}      

\textbf{\textit{na}    \textbf{doktor}}?

to    doctor[\textsc{sg]}

‘How would these elderly women get to a doctor?’ & non-subject PP & Gorni Lom,                                               Belogradchik\\
\lspbottomrule
\end{tabularx}
Inflection class II, feminine and masculine. Nominative and accusative


\begin{tabularx}{\textwidth}{XXXX}

\lsptoprule
11a. & \textbf{\textit{Vod-a}}    \textit{gi}          \textit{e           kara-l-a.}\\
water-\textsc{nom.sg}  \textsc{3pl.dat}   be.\textsc{3sg}   transport-\textsc{ptcp-sg.f}

‘Water moved them.’ & subject NP & Goren Chiflik,                                               Belogradchik\\
11b. & \textit{Sipe-mo} \textbf{\textit{vod-u}}. 

pour- \textsc{prs.1pl}     water-\textsc{acc.sg}

‘We pour water.’ & non-subject NP & Repljana,                                               Belogradchik\\
11c. & \textit{Ide                  i} \textbf{\textit{za} }  \textbf{\textit{vod-u}}.

Go[\textsc{prs.3sg]}   and      for  water-\textsc{acc.sg}

‘And s/he is going to get water.’ & non-subject PP & Gorni Lom,                                               Belogradchik\\
\lspbottomrule
\end{tabularx}
Inflection class III, feminine. Caseless (historically nominative)


\begin{tabularx}{\textwidth}{XXXX}

\lsptoprule
12a. & \textit{…}\textbf{\textit{sol}           }\textit{se        posolva.}

\textsc{…}salt[\textsc{sg}]    \textsc{refl}   salt\textsc{[prs.3sg]}

\textsc{‘}…salt is being added.’ & subject NP & Salash,                                               Belogradchik\\
12b. & \textit{Nasipe-m}    \textbf{\textit{sol} } \textit{i          ga*         pritisne-š}.

pour.in-\textsc{prs.1pl}     salt[\textsc{sg}]   and     \textsc{3sg.acc}   press.down-\textsc{prs.2sg}

‘ I pour in salt and one would press down.’ & non-subject NP & Filipovtsi,                                                                                                              \\
Trân\\
12c. & \textit{Posolva-š          sâs} \textbf{\textit{sol}}.

salt-\textsc{prs.2sg}       with    salt[\textsc{sg}] .

‘You add salt / one adds salt.’ & non-subject PP & Salash,                                               Belogradchik\\
\lspbottomrule
\end{tabularx}
                *Dialectal form corresponding to \textit{go} in Standard Bulgarian.

Inflection class IV, neuter and masculine. Caseless (historically nominative)


\begin{tabularx}{\textwidth}{XXXX}

\lsptoprule
13a. & \textit{…}\textbf{\textit{mlek-o}} \textit{se     podsirva.}

…milk-\textsc{sg}     \textsc{refl}  ferment\textsc{[prs.3sg]}

‘… milk is fermented.’  & subject NP & Chuprene,                                               Belogradchik\\
13b. & \textit{Ako ima-š} \textbf{\textit{mlek-o.}}

if      have-\textsc{prs.2sg}    milk-\textsc{sg}  

‘If you have milk.’ & non-subject NP & Ezdimirtsi,                                                                                                              \\
Trân\\
13c. & \textit{Dve    kofičk-i} \textbf{\textit{ot}    \textbf{mlek-o}}.

two    pot-\textsc{pl}         of     milk-\textsc{sg}  

‘Two pots of milk. ’ & non-subject PP & Ezdimirtsi,                                                                                                              \\
Trân\\
\lspbottomrule
\end{tabularx}
The clear-cut distribution of the nominative and accusative across subject and non-subject positions within inflection class II as presented in \REF{ex:key:11}, however, does not always hold in the border-region dialects. These dialects show a tendency to generalise the nominative for inflection class II nouns across non-subject positions (similar to Standard Bulgarian where this process has been completed in all historical inflection classes, \citealt{Mirčev1958}: 146-147). This results in competing choices for case marking on non-subject NPs and PPs as illustrated in \REF{ex:key:14} and \REF{ex:key:15} by examples from the corpus: 


\begin{tabularx}{\textwidth}{XXXX}

\lsptoprule
14a. & \textbf{\textit{Vod-u}} \textit{gi               dava-mo.} 

water-\textsc{acc.sg}  \textsc{3pl.dat}      give-\textsc{prs.1pl}     

‘We give them water.’ & non-subject NP (accusative) & Repljana,                                               Belogradchik\\
14b. & \textbf{\textit{Vod-a}                }\textit{smo                  kara-l-i.} 

water-\textsc{nom.sg}     be\textsc{[prs.1pl]}      carry-\textsc{ptcp-pl}  

‘We used to carry/haul water.’ & non-subject NP (nominative) & Ezdimirtsi,                                                                                                              \\
Trân\\
\lspbottomrule
\end{tabularx}

\begin{tabularx}{\textwidth}{XXXX}

\lsptoprule
 15a. & \textit{Ne     može} \textbf{\textit{bez}         \textbf{vod-u}               }\textit{da       peče-š}. 

\textsc{neg}   possible  without  water-\textsc{acc.sg}  \textsc{comp}   roast-\textsc{prs.2sg}     

‘One can’t roast [meat] without water.’ & non-subject PP (accusative) & Filipovtsi,                                                                                                              \\
Trân\\
15b. & \textit{Da     ne    ostane} \textbf{\textit{bez}           \textbf{vod-a}}. 

\textsc{comp} \textsc{neg}  remain[\textsc{prs.3sg}]   without   water-\textsc{nom.sg}

‘So that s/he won’t be left without water.’ & non-subject PP (nominative) & Salash,                                               Belogradchik\\
\lspbottomrule
\end{tabularx}
From a diachronic perspective, the situation in Bulgarian border-region dialects reflects one of the final stages in the transition to uninflectedness: only one of four historical inflection classes retains case distinctions, and in only one part of the paradigm (singular). This is different from what we see at the other end of the dialect continuum, the three groups of Serbian dialects in \sectref{sec:key:3}, which illustrate the initial stages in the decomposition of the Proto-Slavic case system. This raises the question as to whether the factors that condition the loss of case distinctions remain stable over time, or whether they change from one stage of the process to the next. In what follows we address this question by looking into the synchronic competition between accusative as a specific non-subject case and nominative as a single form for all syntactic functions, and identify what might trigger speakers’ choices where both alternatives are available (as in inflection class II nouns).

\textbf{4.1.} \textbf{Data} 

The study is based on annotated transcripts of sociolinguistic interviews recorded in West Bulgaria (the dialects of Belogradchik and Trân, part of the border-region dialects). The fieldwork was conducted within the project “Declining case: inflectional loss in progress” by Vladimir Zhobov (University of Sofia, project consultant for Bulgarian) and Alexander Krasovitsky in 2021 and 2022. Additional recordings provided by Vladimir Zhobov stem from his previous fieldwork in North-West Bulgaria. 

 
%%please move the includegraphics inside the {figure} environment
%%\includegraphics[width=\textwidth]{figures/a5Krasovitskyetalrevisedversion-img005.png}

%%[Warning: Draw object ignored]\begin{figure}
\caption{Locations of the interviews in Bulgaria. Border-region dialects} 
\label{fig:key:5}
\end{figure}

%%[Warning: Draw object ignored]
%%[Warning: Draw object ignored]
                         Belogradchik sub-group            Tr $\text{â}$ n sub-group

Forty-six residents born between 1926 and 1950 were recorded (one interview per language consultant; the interviews range in length from forty minutes to two hours). Thirteen of the recorded consultants retain two cases, nominative and accusative, in inflection class II. These thirteen interviews make up a corpus of approximately 168,000 words containing 2292 phrases with inflection class II nouns in non-subject positions. All thirteen interviews demonstrate a competition of accusative and nominative in non-subject positions, however, the probabilities vary considerably between speakers (\figref{fig:key:6}). 

  
%%please move the includegraphics inside the {figure} environment
%%\includegraphics[width=\textwidth]{figures/a5Krasovitskyetalrevisedversion-img006.png}
 

\begin{figure}
\caption{Relative frequency of nominative and accusative forms in inflection class II nouns in non\nobreakdash-subject positions (by speaker).  Corpus size: around 168, 000 tokens. 13 sociolinguistic interviews, 2292 occurrences of inflection class II nouns in non-subject positions. Codes below the columns identify individual speakers.}
\label{fig:key:6}
\end{figure}

\textbf{4.2.} \textbf{Hypothesis:} \textbf{information} \textbf{structure} \textbf{conditions} \textbf{case} \textbf{selection}

Immediate observations made during the sociolinguistic interviews and further elicitation work led us to the following hypothesis. Variation between the nominative and accusative on non-subjects is conditioned by information structure and pragmatic salience: nouns associated with new and pragmatically important information tend to lose specific non-subject case marking and attract nominative; nouns associated with background information and pragmatically less salient parts of an utterance tend to retain specific non-subject marking and show stronger preference for the accusative. 

In what follows, we briefly consider the role of pragmatic factors in case selection within a general linguistic context (\sectref{sec:key:4.3}) and then proceed to the analysis of our data (\sectref{sec:key:4.4}) and a discussion (\sectref{sec:key:4.5}). We argue that pragmatic conditions play a crucial role in speakers’ choices between the two competing cases which mark non-subjects and in the decline of the accusative as distinct morphological marking for non-subjects. At the same time, we will show that the relationship between case and information structure in Bulgarian border-region dialects differs significantly from that observed cross-linguistically, and offer an explanation for this. 

\textbf{4.3.} \textbf{Information} \textbf{structure} \textbf{as} \textbf{a} \textbf{condition} \textbf{on} \textbf{case} \textbf{marking}

Information structure is claimed as a cause for variable case marking cross-linguistically (e.g. \citealt{Aikhenvald2010}, \citealt{Iemmolo2010}, \citealt{Valle2011}; \citealt{DalrympleNikolaeva2011}). First, a large body of typological research links differential object marking (DOM) to disambiguation: typical objects which cannot be confused for subjects are less likely to receive specific grammatical marking than atypical ones, where such marking is essential for disambiguation (\citealt{Silverstein1976}, \citealt{Comrie1977} \textit{inter alia}). With respect to information structure, this would mean that topical objects need to be marked to avoid ambiguity with subjects, while objects occupying their typical slot as part of the pragmatic focus may remain unmarked. In some languages specific non-subject markers are used with highly topical objects dislocated from their default position (\citealt{Iemmolo2010}, \citealt{Aikhenvald2003}).  Second, the degree of pragmatic prominence is another factor which may affect case marking: it has been shown for a number of DOM languages that more prominent objects with a higher pragmatic status are more likely to be marked for case than pragmatically less salient objects (\citealt{Aissen2003}, \citealt{Malchukov2008}). 

It is useful to assess the South Slavic situation against this broad cross-linguistic picture. The data from Bulgarian border-region dialects provide no evidence for case marking being used as a tool in contexts where atypical objects (e.g. topicalised and dislocated to non-default positions or used with omitted governors) may be confused with subjects (see the analysis of the relationship between word order, pragmatic structure and case marking in \sectref{sec:key:4.4}). At the same time, the situation in the border-region dialects is different, in broad terms, to a situation in some of the DOM languages (such as those considered in \citealt{Aissen2003} and \citealt{Malchukov2008}) where pragmatically more salient nouns attract case marking while less pragmatically salient nouns do not. In the Bulgarian dialects in question, case marking on non-subject NPs is conditioned, to a significant extent, on whether they are associated with background or new information in an utterance, or with the topic–focus distinction. However, a significant peculiar feature of these dialects is that specific non-subject case marking (accusative) is more frequent on nouns which are backgrounded. Nominative forms (which may be used in subject and non-subject position, as examples in (14 and \REF{ex:key:15} demonstrate) are more frequent in the pragmatically salient part of an utterance. We will consider this in greater detail in the next section. It should be noted that in qualifying parts of an utterance as topic or focus we follow the tradition of defining topic as existing, background knowledge, or ‘common ground’, and focus as part of an utterance which brings in new information and thus modifies ‘common ground’ (\citealt{HoopSwart2000}; \citealt{Krifka2007}). Since originally our data collection was not intended to test the impact of information structure on case marking, and such analysis was applied later, we used broad definitions based on this opposition, and excluded from the analysis utterances where a non-ambiguous definition of information structure was not possible. Under this view, three types of utterances required special decisions: utterances with contrastive topic as in \REF{ex:key:16}, those with repeated focus, as in \REF{ex:key:17} and those with contrastive focus, as in \REF{ex:key:18}. 

16. Interviewer:   {}- \textit{Kak platiš na toj kojto gledal ovci celoto leto?}\\
                                   ‘ How do you pay someone who looked after sheep the whole summer?’

       Consultant:   \textit{Ne       ima-l-o                          takâv}\textbf{.  \textbf{U}  \textbf{naš-a-ta}                                \textbf{bačij-a}} \\
                              \textsc{neg}  \textsc{exist-ptcp-sg.n}           such.    at   our-\textsc{nom.sg-def.nom.sg}      village.sheep-\textsc{nom.sg}  

 \textit{ne        ima-l-o                         ovčar-e}.                  \\
                              \textsc{neg}  \textsc{exist-ptcp-sg.n}           shepherd-\textsc{pl}               \textbf{\\
}                                                                              \textbf{(contrastive} \textbf{topic)}

       ‘There was no such person. For our village sheep, there were no shepherds.’ 

17.  \textit{A     bab-a                        mi          ode-še…         pase-še} \textbf{\textit{krav-a}}\textit{…}      \\
      and  grandmother-\textsc{nom.sg}    1\textsc{sg}.\textsc{dat}  go-\textsc{impf.3sg}   graze-\textsc{impf.3sg}     cow-\textsc{nom.sg}

 \textit{Ja              njakolko    pât      otiš-l-a           sâs      nju*,}\\
    1\textsc{sg}.\textsc{nom}   several      time    go-\textsc{ptcp-sg.f}    with    \textsc{3sg}.\textsc{acc.f}

 \textit{ama poveče  ona                si         e                        pas-l-a} \textbf{\textit{krav-u}}.\\
    but   more     \textsc{3sg}.\textsc{nom.f}     \textsc{part}    \textsc{aux}.\textsc{prs.3sg}      graze{}-\textsc{ptcp-sg.f}          cow-\textsc{acc.sg}\textbf{\\
}                                                                            \textbf{(repeated} \textbf{focus)}  

‘And my grandmother used to go… grazed the cow. I went with her a few times but mostly she grazed the cow.’\\
*Dialectal form used in the border-region and some other dialects corresponding with the Standard Bulgarian \textit{neja} ‘she [ACC]’ (cf. \citealt{Todorov2002}: 82).                       \\
\\
18. Interviewer:   {}- \textit{A kâde go slagate kopâra?}\\
                                   ‘ Where (i. e. in which dishes) do you add dill?’

       Consultant:   \textit{U tarator tura-m       mnogo}\textbf{.         }\textit{Daže    i}\textbf{   \textbf{u} \textbf{salat-a}               }\textit{si      tura-m.} \\
                              In tarator[\textsc{sg}]  put-1\textsc{sg}  much      even  and  in salad-\textsc{nom.sg}    \textsc{part}   put-1\textsc{sg}  \textbf{\\
}                                                                              \textbf{(contrastive} \textbf{focus)}

      ‘I put a lot of dill in tarator. I add it even in the salad.’

Following the approach taken in this study, we grouped together utterances where non-subject nouns appear in the repeated focus, as in \REF{ex:key:17}, and in the topic on the one hand, and those where non-subject nouns are part of the contrastive topic, as in \REF{ex:key:16}, and of the focus. The rationale behind both decisions is based on information novelty. In the former case, information in the repeated focus is already brought into the discourse and cannot be defined as changing the ‘common ground’. In the latter case, contrastive topic turns out to be very close to focus, in that the contrast itself contributes to existing knowledge and thus modifies the ‘common ground’. The similarity of focus and contrastive topic has been demonstrated in a number of studies within Alternative Semantics, where both concepts are understood as parts of an utterance which provide new information by singling out existing alternatives (\citealt{Rooth1992}: 75-77; \citealt{Büring2016}: 64-68). A small number of examples were annotated as contrastive focus, as in \REF{ex:key:18}. Contrastive focus is not different from focus in that it, like focus, provides new information and modifies the ‘common ground’, under some specific conditions such as polarity of unexpectedness (e. g. \citealt{Zimmermann2008}: 348; \citealt{Goodhue2022}: 117-126). For this reason, in the analysis presented in (4.4) we did not distinguish between the two types of focus. However, a possible effect of contrastive focus will be discussed separately in the account presented in (4.5). 

We will consider the effect of information structure on non-subject case marking with respect to two other factors: word order and phrase type (NP vs. PP).

Previous research has suggested that the diachronic loss of case marking is accompanied by compensatory freezing of word order as a means of distinguishing grammatical roles \citep{Blake2004}. Among other things, this implies that if an object noun occurs outside of its expected position, it will be more likely to preserve its original case marking. This straightforward complementarity between fixed word order and case marking does not necessarily hold (cf. \citealt{Allen2006} for Old \citealt{EnglishDetges2009} for Old French); in fact, the opposite correlation may occur. Data from Japanese show that word order alternations may loosen grammatical relationships within a sentence so that dislocated objects may be less(!) likely to receive expected case marking than their counterparts in neutral positions within a basic constituent structure (cf. Japanese scrambling of accusative-marked objects to a position where only the nominative may be assigned, \citealt{Kasai2018}).

The effect of phrase type (noun phrases vs. prepositional phrases) on case selection and decline of one of the competing cases is surveyed in \sectref{sec:key:2}. Here we will only reiterate the observation made above that the correlation between the use of prepositions and the tendency to maintain an original case form has been proved to be strong for the dialect types under study. Therefore, we found it important to disentangle different conditions on case marking in order to assess their effect individually without interference from other conditions. To do so, we will analyse the effect of topic-focus distinction separately for different types of word order and different phrase types (NP and PP). 

\textbf{4.4.} \textbf{The} \textbf{effect} \textbf{of} \textbf{topic} \textbf{and} \textbf{focus} \textbf{on} \textbf{case} \textbf{marking} \textbf{in} \textbf{Bulgarian} \textbf{border-region} \textbf{dialects}

To investigate the effect of information structure on competing choices in case marking, a sub-corpus of approximately 90,000 words (six sociolinguistic interviews) was annotated with respect to information structure. The percentage of nominative on non-subjects in the six included interviews is as follows: 7\% (BEL1-12), 20\% (BEL1-15), 23\% (BEL1-31), 42\% (KS1-09), 55\% (KS1-08), 57\% (KS1-18).  A dataset of 956 utterances containing inflection class II nouns in non-subject positions and allowing unambiguous assignment of the five values (topic, contrastive topic, focus, contrastive focus, repeated focus) was extracted. The data were combined for (i) topic and repeated focus, and for (ii) focus, contrastive focus and contrastive topic, as discussed in \sectref{sec:key:4.3}. The analysis presented in this section is based on these two subsets. 

  
%%please move the includegraphics inside the {figure} environment
%%\includegraphics[width=\textwidth]{figures/a5Krasovitskyetalrevisedversion-img007.png}
 

\begin{figure}
\caption{Nominative and accusative on non-subject nouns} 
\label{fig:key:7}
\end{figure}

%%[Warning: Draw object ignored]
%%[Warning: Draw object ignored]
%%[Warning: Draw object ignored]
%%[Warning: Draw object ignored]
As \figref{fig:key:7} demonstrates, there is a significant difference in the frequency of the nominative and accusative across the two subsets. The accusative dominates in topic \REF{ex:key:19a} or repeated focus position \REF{ex:key:19b}: 79\% accusative forms. Frequency of the accusative falls significantly in focus \REF{ex:key:19c} and contrastive topic positions \REF{ex:key:19d}: 49\%. In other words, non-subject nouns are more likely to take the accusative in pragmatically less salient positions; pragmatically more salient parts of a sentence show a stronger preference for the nominative on non-subjects.

19a.  \textit{I       slam-a-ta}   \textit{otzad.    i       otvârlj-u} \textbf{\textit{slam-u-tu}}   \\
        and  straw-\textsc{nom.sg-def.nom.sg}    behind  and   toss-\textsc{prs.3pl} straw-\textsc{acc.sg-def.acc.sg}\\
 \textit{nastranu.\\
}         aside                                                                                                  \textbf{(topic)}

‘And the straw is behind. And they toss the straw aside.’

19b. \textit{Čorapete  ot} \textbf{\textit{dvojk-a}} \textit{si   plete-mo}. \\
          socks  from    double.thread-\textsc{nom.sg}  \textsc{part}  weave- \textsc{prs.}1\textsc{pl}\\
                                         \textbf{(focus)}

         \textit{I}    \textit{ot}   \textbf{\textit{dvojk-u} }               \textit{i     ot}    \textit{trojk-u}                       \textit{sam               ple-l-a.}                 \\
        and from  double.thread-\textsc{acc.sg} and from   triple.thread-\textsc{acc.sg} be\textsc{.prs.}1\textsc{sg}] weave-\textsc{ptcp-sg.f}\\
                        \textbf{(repeated} \textbf{focus)}\\
       ‘We weave socks with double threads. I have weaved both with double and triple threads.’

19c.  \textit{Ako   nema         rek-a}                  \textit{nosi-mo} \textbf{\textit{vod-a.}}\\
         if      \textsc{exist.neg}    river-\textsc{nom.sg}  carry- \textsc{prs.}1\textsc{pl}     water-\textsc{nom.sg}\\
                                                                                 \textbf{(focus)}\\
         ‘If there is no river, we carry water.’

19d.  Interviewer: \textit{S râka li pravete?}\\
                                  ‘Do you do it manually (“with a hand”)?’

           Consultant: \textit{S       râk-u,              a        s       kvo?} \\
                               with  hand-\textsc{acc.sg}   and    with what

                               \textbf{\textit{S}       \textbf{lev-a-ta}} \textit{tegli-š              otgore,} \\
                               with left-\textsc{nom.f.sg-def.nom.sg} \textsc{f}  pull-\textsc{prs.2sg}    from.above                                               \\
                                \textbf{\textit{s}        \textbf{desn-a-ta}} \textit{vârti-š                 vreten-o-to}\\
                                with right-\textsc{nom.f.sg-def.nom.sg.f}   rotate-\textsc{prs.2sg}   spindle-\textsc{sg-def.sg.n}\\
                                      \textbf{(contrastive} \textbf{topic)}

 ‘Manually, how else? With your left hand you pull it from above, with your right hand you rotate the spindle.’ 

We assessed the effect of information structure on case marking in connection with two potentially relevant factors, phrase type (presence or absence of a preposition) and word order. As shown above, phrase type proved to have an effect on case marking in the Serbian dialects at an earlier stage of case decline (\sectref{sec:key:3}). Word order, although not independent of information structure, may still vary within the same topic–focus structure, therefore this factor was also disentangled from the other two factors under study. In Bulgarian, direct and indirect objects under basic default word order occur postverbally. Non-default ordering (e.g. preverbal objects) is motivated by a variety of conditions, which may or may not be related to information structure (\citealt{Georgieva1974}; \citealt{Georgieva1983}; \citealt{Dyer1992}). Thus, topicalised objects are normally fronted, but this is not a universal rule, and under appropriate contextual conditions topicalised objects may occur postverbally \REF{ex:key:20a}. 

20a.  Consultant 1: \textit{Kâde     utiva-š?}\\
                                   Where  go-\textsc{prs.2sg}?\\
                                           

         Consultant 2: \textit{Po rabot-a,             majk-o}.\\
                                  for work-\textsc{nom.sg}   mother-\textsc{voc}

          Consultant 1: \textit{Ah, mrazj-a             az} \textbf{\textit{tazi}                    \textbf{rabot-a}}!\\
                                   ah   hate-\textsc{prs.1sg}   \textsc{1sg}.\textsc{nom}  this[\textsc{nom.sg.f}]  work-\textsc{nom.sg}

‘C1: Where are you going? C2: On business, mother. C1: Oh I hate this business!’

On the other hand, as expected in SVO languages, objects in focus normally occur post-verbally, but strong logical stress may result in emphatic word order, under which an object occurs in preverbal position \REF{ex:key:20b}. 

20b.  \textbf{Cjal-a-ta}                                        \textbf{planin-a}                      vdišva-m \\
  whole- \textsc{nom.sg.f-def.nom.sg.f}  mountain- \textsc{nom.sg}     inhale-\textsc{prs.1sg}

‘I inhale all this mountain range’ (adapted from \citealt{Georgieva1983}: 284). 

In addition to postverbal and preverbal ordering, we found a significant number of elliptic constructions, for example, objects with a missing governing verb, or interruptions where a governor and a governee stand far apart from each other in different syntagms. To see whether the absence of an overt governor has an effect on case marking, such phrases (tagged as ‘isolated’) were considered separately. 

Therefore, in order to disentangle the contribution of potentially conditioning factors, we assessed the effect of information structure separately for six different combinations of word order and phrase type: 

21a.     NP, postverbal

\textit{Ne     smo              slaga-l-i} \textbf{\textit{vod-u}}\\
\textsc{neg}    \textsc{aux.prs}.1\textsc{pl}  put-\textsc{ptcp-pl}    water-\textsc{acc.sg}\\
‘We added no water.’

21b.     NP, preverbal

\textbf{\textit{Krav-a} } \textit{smo                 ima-l-i}\\
cow-\textsc{nom.sg}  \textsc{aux.prs}.1\textsc{pl}     have-\textsc{ptcp-pl} \\
‘We had a cow.’

21c.     NP, isolated

\textit{Ne    sâm       ima-l-a                bab-u,}\\
\textsc{neg}   \textsc{aux}.1\textsc{sg} have-\textsc{ptcp-sg.f}   grandmother-\textsc{acc.sg}

\textit{samo} \textbf{\textit{majk-a}}\\
only    mother-\textsc{nom.sg}\\
‘I didn’t have a grandmother, just a mother.’

21d.    PP, postverbal

\textit{I        ojde-mo        večerom} \textbf{\textit{na}}  \textbf{\textit{večerj-u}}\\
and  go-\textsc{prs.1pl}     in.the.evening    to  evening.meal-\textsc{acc.sg}\\
‘And in the evening we will go to have evening meal.’

21e.     PP, preverbal

\textbf{\textit{Na} \textbf{bašt-a}                 }\textit{mi           pomaga-x.}\\
to  father-\textsc{nom.sg}   \textsc{1sg.dat}   help-\textsc{aor.1sg}\\
‘I helped my father.’  

21f.      PP, isolated

Interviewer: \textit{– Ima li češma?}  \\
                            ‘Is there a tap?’

Consultant: \textit{– Otzad             češm-a-ta.} \textbf{\textit{Na} \textbf{kuxničk-u.}}\\
                         in.the.rear    tap-\textsc{nom.sg-def.nom.sg}      in   kitchen-\textsc{acc.sg} \\
‘The tap is in the rear. In the kitchen.’

Relative frequencies of the two alternative cases, nominative and accusative, on non-subject nouns with respect to the six combinations of conditions are presented in \figref{fig:key:8} (for NPs) and \figref{fig:key:9} (for PPs). 


\begin{tabularx}{\textwidth}{XXX}

\lsptoprule
 POSTVERBAL NP & PREVERBAL NP & ISOLATED NP\\
  
%%please move the includegraphics inside the {figure} environment
%%\includegraphics[width=\textwidth]{figures/a5Krasovitskyetalrevisedversion-img008.png} & %%please move the includegraphics inside the {figure} environment
%%\includegraphics[width=\textwidth]{figures/a5Krasovitskyetalrevisedversion-img009.png} & %%please move the includegraphics inside the {figure} environment
%%\includegraphics[width=\textwidth]{figures/a5Krasovitskyetalrevisedversion-img010.png}
 \\
& %%[Warning: Draw object ignored]
%%[Warning: Draw object ignored]
%%[Warning: Draw object ignored]
%%[Warning: Draw object ignored] & \\
\lspbottomrule
\end{tabularx}
\begin{figure}
\caption{Preverbal, postverbal and isolated NPs in topic and focus positions}
\label{fig:key:8}
\end{figure}


\begin{tabularx}{\textwidth}{XXX}

\lsptoprule
 POSTVERBAL PP & PREVERBAL PP & ISOLATED PP\\
  
%%please move the includegraphics inside the {figure} environment
%%\includegraphics[width=\textwidth]{figures/a5Krasovitskyetalrevisedversion-img011.png} & %%[Warning: Draw object ignored]
%%[Warning: Draw object ignored]
%%[Warning: Draw object ignored]
  
%%please move the includegraphics inside the {figure} environment
%%\includegraphics[width=\textwidth]{figures/a5Krasovitskyetalrevisedversion-img012.png} & %%please move the includegraphics inside the {figure} environment
%%\includegraphics[width=\textwidth]{figures/a5Krasovitskyetalrevisedversion-img013.png}
 \\
\multicolumn{3}{c}{%%[Warning: Draw object ignored]
}\\
\lspbottomrule
\end{tabularx}
\begin{figure}
\caption{Preverbal, postverbal and isolated PPs in topic and focus positions} 
\label{fig:key:9}
\end{figure}

A key outcome of this analysis is the uniform pattern which we found across all six data subsets (factor combinations): while pragmatically more important parts of an utterance (focus or contrastive topic) demonstrate a strong tendency to generalize the nominative on non-subject nouns, the accusative as a distinct non-subject case is better preserved in pragmatically less salient positions (topic and repeated focus). In the light of the findings reported above for Serbian dialects, it is particularly striking that structural conditions, that is, the absence or presence of a preposition, seem to have no impact on speakers’ choices, and neither enhance the use of a distinct non-subject case form (accusative), nor contribute to its loss and the generalisation of an undifferentiated morphological form (nominative). In other words, non-subject NPs and PPs in terms of their case marking show similar sensitivity to information structure. With respect to word order, a strong effect of information structure is found for all three linear structures: postverbal, preverbal and isolated positions. In all three data subsets, the accusative is better retained in topic and repeated focus position, and is less frequent in focus and contrastive topic position where it is replaced by the nominative. 

The statistical significance of the results presented in \figref{fig:key:8} and \figref{fig:key:9} was assessed by ANOVA analysis. The effect of each of the three factors, phrase type, information structure and word order, was considered as a sole factor, and in connection with the other two factors (the null hypothesis assumed that there was no effect of either of these factors on case selection). The outcome of this analysis is presented in \tabref{tab:key:4}. 


\begin{tabularx}{\textwidth}{XXX}

\lsptoprule

\textbf{Phrase} \textbf{type} \textbf{(NP} \textbf{vs.} \textbf{PP)} & \textbf{Word} \textbf{order}  & \textbf{Information} \textbf{structure}\\
NO EFFECT

p = 0.35 in connection with WO and IS

p = 0.29 taken separately & WEAK EFFECT (?)

p = 0.18 in connection with WO and phrase type  

p < 0.001 taken separately & STRONG EFFECT

p < 0.001 under any condition(s) or taken separately\\
\lspbottomrule
\end{tabularx}
%%please move \begin{table} just above \begin{tabular
\begin{table}
\caption{The effect of phrase type, word order and information structure on case marking in Bulgarian border-region dialects}
\label{tab:key:4}
\end{table}

The results of the statistical analysis presented in \tabref{tab:key:4} confirm the observation that phrase type (NP or PP), whether taken separately or in connection with word order and information structure, has no effect on case marking on non-subjects (p=0.29 and 0.35 accordingly). On the other hand, we see a strong effect of information structure, either considered on its own (as in \figref{fig:key:7}), or in connection with phrase type and word order, for sub-sets presented in Figures 8 and 9 (p < 0.001 either taken separately or in connection with other conditions). We cannot, however, arrive at such an unequivocal conclusion on the role of word order. Within each of the six sub-sets presented in Figures 8 and 9, word order does not reach clear statistical significance (p=0.18). Taken as a sole factor without considering other factors (structure of the phrase and its affiliation with different pragmatic roles) postverbal positions seem to enhance the use of the accusative, and pre-verbal positions favour the nominative (p < 0.001). However, it should be noted that the latter outcome may reflect a natural bias in the data. Non-subject postverbal focus NPs are more frequent than preverbal ones. Given that focus tends to attract nominative marking according to our data, this may significantly bias the analysis of the role of word order as a sole factor in case marking. Therefore, we think it important that word order, when taken as a factor in combination with information structure, fails to reach unambiguous statistically significant values. Firm conclusions, however, could only be reached from data normalised with respect to different types of word order and different pragmatic roles, which would require a larger corpus than we currently have. 

\textbf{4.5.} \textbf{Discussion:} \textbf{pragmatic} \textbf{factors} \textbf{in} \textbf{case} \textbf{decline}

The reported findings leave little doubt that information structure has a crucial effect on variation in case marking in Bulgarian border-region dialects. What might this synchronic variation tell us about the diachrony of case decline? We hypothesise that the historical competition of case marking here is motivated by pragmatic salience. What we observe, however, is not a tight coupling between a pragmatic role and grammatical form, but rather a tendency for caseless forms (historically the nominative) to be generalised first across the pragmatically most salient positions, and only then to spread further to other parts of the utterance. Key evidence supporting this view comes from the comparison of different parts of the subcorpus that have been annotated for information structure. As indicated in \sectref{sec:key:4.4} and illustrated by \figref{fig:key:6}, preferences for case marking on non-subject nouns in this subcorpus vary enormously across speakers, ranging from a total of 7\% to 57\% of nominative forms. We divided the subcorpus into two parts depending on the by-speaker frequency of nominative forms. The first part includes the three most conservative speakers (from 7\% to 23\% of nominative forms in non-subject positions). The second part includes interviews recorded from more innovative speakers with an overall frequency of nominative forms between 42\% and 57\% (T-test confirms the significance of the difference between the two groups with respect to case variation: p<0.0001, t=8.4). 

The analysis presented below shows which parts of the information structure are most susceptible to the loss of case distinctions at a very early stage of this process, to identify the role of pragmatic factors. We looked at the data from the three most conservative speakers in order to identify conditions which make these speakers deviate from their dominant pattern, accusative marking of non-subjects. We assessed the impact of each pragmatic condition for which our data are annotated. Unlike the findings from the whole subcorpus analysed in the previous section (six interviews), our results here show that in the three most conservative individual dialects, the accusative is strongly associated not only with the background or repeated information (topic and repeated focus), but also with new information (focus). If, however, we look separately at the pragmatically most salient NPs and PPs, contrastive topic and contrastive focus, we find that these pragmatic conditions favour the nominative (\tabref{tab:key:5}, illustrated by examples in \REF{ex:key:22}):


\begin{tabularx}{\textwidth}{XXXXX}

\lsptoprule
%\hhline%%replace by cmidrule{----~}
~ & ACCUSATIVE & NOMINATIVE & \% NOMINATIVE & \\
%\hhline%%replace by cmidrule{----~} &  &  &  & \\
%\hhline%%replace by cmidrule{----~}
Topic & 17 & 0 & 0\% & \\
%\hhline%%replace by cmidrule{----~}
Repeated focus & 13 & 4 & 24\% & \\
%\hhline%%replace by cmidrule{----~}
Focus & 104 & 18 & 15\% & \\
%\hhline%%replace by cmidrule{----~}
Contrastive topic & 0 & 3 & 100\% & \\
%\hhline%%replace by cmidrule{----~}
Contrastive focus & 0 & 2 & 100\% & \\
%\hhline%%replace by cmidrule{----~}
\textbf{Total}  & \textbf{134} & \textbf{27} & \textbf{17\%} & \\
%\hhline%%replace by cmidrule{----~}
\lspbottomrule
\end{tabularx}
%%please move \begin{table} just above \begin{tabular
\begin{table}
\caption{Distribution of nominative and accusative forms with non-subject nouns for the three most conservative individual dialects. Raw numbers and \% of nominative}
\label{tab:key:5}
\end{table}

22a.  Interviewer: - \textit{Kak platiš na toj kojto gledal ovci celoto leto?}\\
    How do you pay someone who looked after sheep the whole summer? 

 Consultant:   \textit{Ne       ima-l-o                  takâv}\textbf{. \textbf{U}  \textbf{naš-a-ta}                                  \textbf{bačij-a}} \\
                                         \textsc{neg}      \textsc{exist-ptcp-sg.n}     such.  at our-\textsc{nom.sg.f-def.nom.sg.f}     village.sheep-\textsc{nom.sg}   

\textit{ne        ima-l-o                 ovčar-e}.                                            \\
\textsc{neg}       \textsc{exist-ptcp-sg.n}     shepherd-\textsc{pl}

                                                                                        \textbf{(contrastive} \textbf{topic)}

       ‘There was no such person. For our village sheep, there were no shepherds.’

22b.  Interviewer: [- \textit{Da mu platiš posle?} \\
 Do you pay him afterwards?

 Consultant:   \textit{{}-  Ne, da ne mu platiš.   \\
} {}- No, you do not pay him.]      

 \textit{Odi-l-i             smo} \textbf{\textit{na}   \textbf{zared-a}}.\\
                                       go-\textsc{ptcp-pl}     aux.\textsc{prs.1pl}   on  work.for.smb-\textsc{nom.sg} \\
                                                                          \textbf{(contrastive} \textbf{focus)}\\
             ‘We went to work for him (in exchange for his favour).’    

Under non-contrastive focus we can expect the accusative (14 c): 

22c.    \textit{i        posle ide-mo} \textbf{\textit{za}  \textbf{nevest-u{}-tu}}.\\
 and after   go-\textsc{prs.1pl}   for bride-\textsc{acc.sg-def.acc.sg} \\
                                              \textbf{(non-contrastive} \textbf{focus,} \textbf{listing} \textbf{consecutive} \textbf{actions)}

            ‘And then we go to pick up the bride.’  

The data from the three most innovative dialects reveal the spread of nominative from the most salient pragmatic conditions, such as contrastive topic and contrastive focus, to the parts of the utterance containing new information in general (as demonstrated in \tabref{tab:key:6}). 


\begin{tabularx}{\textwidth}{XXXXX}

\lsptoprule
%\hhline%%replace by cmidrule{----~}
~ & ACCUSATIVE & NOMINATIVE & \% NOMINATIVE & \\
%\hhline%%replace by cmidrule{----~} &  &  &  & \\
%\hhline%%replace by cmidrule{----~}
Topic & 70 & 33 & 32\% & \\
%\hhline%%replace by cmidrule{----~}
Repeated focus & 69 & 10 & 13\% & \\
%\hhline%%replace by cmidrule{----~}
Focus & 272 & 326 & 55\% & \\
%\hhline%%replace by cmidrule{----~}
Contrastive topic & 1 & 4 & 80\% & \\
%\hhline%%replace by cmidrule{----~}
Contrastive focus & 1 & 9 & 90\% & \\
%\hhline%%replace by cmidrule{----~}
\textbf{Total}  & \textbf{413} & \textbf{382} & \textbf{48\%} & \\
%\hhline%%replace by cmidrule{----~}
\lspbottomrule
\end{tabularx}
%%please move \begin{table} just above \begin{tabular
\begin{table}
\caption{Distribution of nominative and accusative forms with non-subject nouns for the three most innovative individual dialects. Raw numbers and \% of nominative}
\label{tab:key:6}
\end{table}

The data for a subset of the most conservative speakers in the corpus, however small, may be taken as an indication of a diachronic tendency when compared to the whole corpus. While accusative dominates in the speech of the three conservative speakers overall, the pragmatically most salient parts of an utterance favour the nominative on non-subject nouns. This process has further developed with more innovative speakers with whom nominative has spread further into the focus position, where it dominates, and into topic position, though here it is still in the minority. 

The fact that pragmatics play a crucial role in speakers’ choices in Bulgarian border-region dialects may be due to both language-internal and external factors. As we noted earlier in \sectref{sec:key:4}, inflection class~II nouns are the only domain where case distinctions are still preserved, while nouns of all other classes (and adjectival forms which agree with them) are uninflected. This contributes to the marginalisation of case and restricts accusative to less salient positions. This process is supported by external factors, namely exposure to language varieties without case distinctions on a daily basis (younger speakers, TV, written language).This can marginalise accusative forms, restricting them to less highlighted parts of an utterance.  

\textbf{5.} \textbf{Conclusion} 

\begin{stylexelementtoproof}
The ongoing loss of case in Serbian and Bulgarian dialects results in the competition between two forms, one of them more specific and one of them more general. How exactly this competition is manifested depends on the particular case system it occurs in. In more conservative varieties that retain in principle a six-case system (Serbian), competition is between a more specific case (e.g. the genitive) and a general oblique (accusative). In more innovative varieties (Bulgarian), where case loss is more advanced, the competition is between the general oblique form (accusative) and an uninflected stem. A particularly interesting result of our study has been to uncover key factors underlying these systems of competing case marking. Broadly construed, they match what has been observed cross-linguistically, but the principles that govern them are different.  Elsewhere, prior work has shown that in case systems which allow for more explicit and less explicit marking for a given syntactic role, more explicit case marking has two major functions: (i) indicating the syntactic role of nominals in contexts where this is not otherwise obvious, and (ii) highlighting nominals that are semantically or pragmatically prominent. In the Serbian and Bulgarian dialects investigated here, the factors conditioning the competition of cases appear to be the exact opposite. 
\end{stylexelementtoproof}

\begin{stylexelementtoproof}
First, contexts where the syntactic role of a nominal is unambiguously signalled, as in prepositional phrases, in fact provide a more favourable environment for the preservation of the original case forms. We have illustrated this with the data from the Serbian dialects, where the original genitive case forms alternate with the innovative accusative forms. We analysed three construction types: prepositional constructions, adnominal constructions, and constructions with a quantifier. In all the dialects of the sample, prepositional constructions preserved the original case forms better than the other two. The three constructions vary with respect to the strength of connection between their elements. While the use of a noun (such as 'cup' or 'top') or a quantifier does not necessarily entail the subsequent use of a noun, a preposition obligatorily governs a noun phrase, marked with the genitive in our case. Thus, the established association of a form with its apparent context turns out to be a key factor which determines speakers’ choices and the contributes to the preservation of the original case forms. 
\end{stylexelementtoproof}

Second, higher pragmatic salience, rather than triggering case distinctions, as has been widely observed cross-linguistically, in fact leads to the opposite result in the Bulgarian border-region dialects. These dialects retain at most two cases, nominative and accusative. The accusative is normally used for non-subjects, but this is not always so: they may also take the nominative, making them morphologically indistinguishable from subjects and thus effectively uninflected. This variation is conditioned by information structure: a subject (nominative) vs. non-subject (accusative) distinction is better maintained where background (given) information is being presented, while pragmatically salient parts of an utterance show a strong tendency to eliminate this distinction. Additional factors that play a role in limiting the use of the accusative are the fact that only one class of nouns is affected (inflection class II), and the overall marginalisation of case distinctions on nouns in the context of the prestige standard language, which lacks nominal case.  

The fine-grained analysis of the competing patterns of case marking provided here contributes to our understanding of factors that condition the loss of inflectional case. In particular, different factors appear to be activated at different stages of the change. A synchronic cross-dialect analysis therefore allows us to uncover the intricate detail of this historical process.
\begin{verbatim}%%move bib entries to  localbibliography.bib
@book{Aikhenvald2003,
	address = {Cambridge},
	author = {Aikhenvald, Alexandra J.},
	publisher = {Cambridge University Press},
	title = {\textit{A grammar of {Tariana}}},
	year = {2003}
}


@book{Aikhenvald2010,
	address = {Oxford},
	author = {Aikhenvald, Alexandra},
	publisher = {Oxford University Press},
	title = {\textit{Language contact in Amazonia}},
	year = {2010}
}


@book{Aissen2003,
	address = {economy. \textit{Natural Language \& Linguistic Theory} 21 \REF{ex:key},
	author = {Aissen, Judith},
	publisher = {3}, 435-483},
	title = {Differential object marking: {{I}}conicity vs},
	year = {2003}
}


@incollection{Allen2006,
	address = {Oxford},
	author = {Allen, Cynthia},
	booktitle = {\textit{The Handbook of the History of {English}}},
	editor = {A. van Kemenade and B. Los},
	pages = {201–223},
	publisher = {Blackwell},
	title = {Case Syncretism and Word Order Change},
	year = {2006}
}


@book{Belić1905,
	address = {Srpski dijalektološki zbornik \REF{ex:key:1}.} Belgrade},
	author = {Belić, Aleksandr},
	publisher = {Državna Štamparija Kraljevine Srbije},
	sortname = {Belic, Aleksandr},
	title = {\textit{Dijalekti Istočne i južne Srbije},
	year = {1905}
}


@book{Belić1999,
	address = {\textit{Istorija srpskog jezika: fonetika, reči sa deklinacijom, reči sa konjugacijom.} Belgrade},
	author = {Belić, Aleksandar},
	publisher = {Zavod za udžbenike i nastavna sredstva},
	sortname = {Belic, Aleksandar},
	year = {1999}
}


@book{Bernštejn1948,
	address = {Vol. 1. Jazyk valašskix gramot XIV-XV vekov}. Moscow – Leningrad},
	author = {Bernštejn, Samuil B.},
	publisher = {AN SSSR},
	sortname = {Bernstejn, Samuil B.},
	title = {\textit{Razyskanija v oblasti bolgarskoj istoričeskoj dialektologii},
	year = {1948}
}


@book{Blake2004,
	address = {\textit{Case.} (2\textsuperscript{nd} edition) Cambridge},
	author = {Blake, Barry J.},
	publisher = {Cambridge University Press},
	year = {2004}
}


\textstyleauthor{Bossong, Georg.}\textstylepubyear{1991}.~\textstylechaptertitle{Differential object marking in Romance and beyond}. In~\textstyleeditor{D. Kibbee}~\& ~\textstyleeditor{D. Wanner}~(eds.).~\textstylebooktitle{\textit{New~analyses in Romance~linguistics}}. \textstylepagefirst{143}{}-\textstylepagelast{170}.~\textstylepublisherlocation{Amsterdam}: John Benjamins.

Büring, Daniel. 2016. (Contrastive) Topic. In Büring, Daniel, Féry, Caroline \& Ishihara, Shinichiro (eds.). The \textit{Oxford Handbook of Information Structure,} 64-85. Oxford: Oxford University Press.

@book{Češko1970,
	address = {Moscow},
	author = {Češko, Elena B.},
	publisher = {Nauka},
	sortname = {Cesko, Elena B.},
	title = {\textit{Istorija bolgarskogo sklonenija}},
	year = {1970}
}


@article{Comrie1977,
	author = {Comrie, Bernard},
	journal = {\textit{Études Finno-Ougriennes}},
	note = {. Budapest: Akadémiai Kiadó.},
	pages = {5--17},
	title = {Subjects and direct objects in {Uralic} languages: {{A}} functional explanation of case-marking systems},
	volume = {12},
	year = {1977}
}


@book{Comrie1989,
	address = {Chicago},
	author = {Comrie, Bernard},
	publisher = {University of Chicago Press},
	title = {\textit{Language universals and linguistic typology}},
	year = {1989}
}


@incollection{Cristofaro2013,
	address = {Berlin},
	author = {Cristofaro, S.},
	booktitle = {\textit{Languages across boundaries: {{S}}tudies in memory of {Anna} {Siewierska}} },
	editor = {D. Bakker, and M. Haspelmath},
	pages = {69--93},
	publisher = {de Gruyter Mouton},
	title = {The referential hierarchy: {{R}}eviewing the evidence in diachronic perspective},
	year = {2013}
}


@book{DalrympleDalrymple2011,
	address = {Cambridge},
	author = {Dalrymple, Mary and Irina Nikolaeva},
	publisher = {Cambridge University Press},
	title = {\textit{Object and Information Structure}},
	year = {2011}
}


Detges, Ulrich. 2009. How useful is case morphology? The loss of the Old French two-case system within a theory of preferred argument structure. In J. Barðdal \& S. L. Chelliah (eds). \textit{The role of semantic, pragmatic, and discourse factors in the development of case,} 93-120. Amsterdam: John Benjamins.

@book{Duridanov1956,
	address = {Godišnik na Sofijskija universitet, 51 \REF{ex:key:1}. Sofia},
	author = {Duridanov, Ivan},
	publisher = {Nauka i izkustvo},
	title = {\textit{Kâm problemata za razvoja na bâlgarskija ezik ot sintetizâm kâm analitizâm}},
	year = {1956}
}


@book{Duridanov1958,
	address = {\textit{Rocznik slawistyczny} XX \REF{ex:key},
	author = {Duridanov, Ivan},
	note = {16-26.},
	publisher = {2}},
	title = {Za načenkite na analitizma v bâlgarskija ezik},
	year = {1958}
}


@book{Dyer1992,
	address = {Amsterdam},
	author = {Dyer, Donald L.},
	publisher = {Rodopi},
	title = {\textit{Word order in the simple {Bulgarian} sentence: {{A}} study in grammar, semantics and pragmatics}},
	year = {1992}
}


@book{Feleško1955,
	address = {\textit{Značenje i sintaksa srpskohrvatskog genitiva.} Belgrade: Orfelin – Belgrade: Vukova zadužbina – Novi Sad},
	author = {Feleško, Kazimjež},
	publisher = {Matica Srpska},
	sortname = {Felesko, Kazimjez},
	year = {1955}
}


@book{Georgieva1974,
	address = {Sofia},
	author = {Georgieva, Elena},
	publisher = {Bâlgarska akademija na naukite},
	title = {\textit{Slovored na prostoto izrečenie v bâlgarskija knižoven ezik}},
	year = {1974}
}


@incollection{Georgieva1983,
	address = {Sofia},
	author = {Georgieva, Elena},
	booktitle = {\textit{Gramatika na sâvremennija bâlgarski knižoven ezik}. {{V}}ol. 3},
	editor = {K. Popov},
	pages = {266--271},
	publisher = {Bâlgarska akademija na naukite},
	title = {Slovored},
	year = {1983}
}


@article{Goodhue2022,
	author = {Goodhue, Daniel},
	journal = {\textit{Journal of Semantics}},
	number = {1},
	pages = {117--158},
	title = {All Focus is Contrastive: {{O}}n Polarity (Verum) Focus, Answer Focus, Contrastive Focus and Givenness},
	volume = {39},
	year = {2022}
}


@book{HewsonHewson2006,
	address = {The Development of Configurational Syntax in Indo-European Languages.} Amsterdam},
	author = {Hewson, John and Vit Bubenik},
	publisher = {John Benjamins},
	title = {\textit{From Case to Adposition},
	year = {2006}
}


@book{Hjelmslev1935,
	address = {Étude de grammaire générale}. Aarhus},
	author = {Hjelmslev, Louis},
	publisher = {Universitetsforlaget},
	title = {\textit{{La} catégorie des cas},
	year = {1935}
}


Hoop, Helen de \& Henriette de Swart. 2000. Topic and focus. In Lisa Cheng \& Rint Sybesma (eds.). \textit{The First Glot International State-of-the-Article Book. The Latest in Linguistics,} 105-130. Berlin: De Gruyter Mouton.

@article{Iemmolo2010,
	author = {Iemmolo, Giorgio},
	journal = {\textit{Studies in Language}},
	number = {2},
	pages = {239--272},
	title = {Topicality and differential object marking: {{{E}}}vidence from {Romance} and beyond},
	volume = {34},
	year = {2010}
}


@book{Ivić1956,
	address = {\textit{Dijalektologija srpskohrvatskog jezika: Uvod u štokavsko narečje.} Novi Sad},
	author = {Ivić, Pavle},
	publisher = {Matica Srpska},
	sortname = {Ivic, Pavle},
	year = {1956}
}


@book{Kasai2018,
	address = {\textit{Glossa: a journal of general linguistics} 3(1): 127. doi},
	author = {Kasai, Hironobu},
	note = {org/10.5334/gjgl.676}.},
	publisher = {\url{https://doi},
	title = {Case valuation after scrambling: {{N}}ominative objects in {Japanese}},
	year = {2018}
}


\textstyleauthor{König, Christa.}~\textstylepubyear{2008}.~\textstylebooktitle{\textit{Case in Africa}}. Oxford: Oxford University Press.

@incollection{Krifka2007,
	address = {Potsdam},
	author = {Krifka, Manfred},
	booktitle = {\textit{The notions of information structure}},
	editor = {Caroline Féry, Gisberg Fanselow and Manfred Krifka},
	pages = {13--55},
	publisher = {Universitätsverlag},
	title = {Basic notions of information structure},
	year = {2007}
}


@incollection{Kulikov2006,
	address = {Amsterdam},
	author = {Kulikov, Leonid},
	booktitle = {\textit{Case, Valency and Transitivity}},
	editor = {Leonid Kulikov, Andrej Malchikov and Peter de Swart},
	pages = {23--49},
	publisher = {John Benjamins},
	title = {Case systems in a diachronic perspective: {{A}} typological sketch},
	year = {2006}
}


\textstyleauthor{Kurumada,} Chigusa \& ~Tim Florian \textstyleauthor{Jaeger.}~\textstylepubyear{2015}.~\textstylearticletitle{Communicative efficiency in language production: Optional case-marking in Japanese}.~\textit{Journal of Memory and Language}~\textstylevol{83}.~\textstylepagefirst{152}{}-\textstylepagelast{178}.

@book{Malchukov2008,
	address = {\textit{Lingua} 118 \REF{ex:key},
	author = {Malchukov, Andrej},
	note = {203-221.},
	publisher = {2}},
	title = {Animacy and asymmetries in differential case marking},
	year = {2008}
}


@incollection{McGregor2018,
	address = {Berlin: Language Science Press. DOI},
	author = {McGregor, William},
	booktitle = {\textit{Diachrony of differential argument marking. (Studies in Diversity Linguistics 19)}},
	editor = {Ilja A. Seržant and Alena Witzlack-Makarevich},
	note = {5281/zenodo.1219168},
	pages = {243--279},
	publisher = {10},
	title = {Emergence of optional accusative case marking in {Khoe} languages},
	year = {2018}
}


@book{Meje2001,
	address = {Moscow},
	author = {Meje, Antuan},
	publisher = {Progress},
	title = {\textit{Obščeslavjanskij jazyk}},
	year = {2001}
}


@book{Miloradović2003,
	address = {\textit{Upotreba padežnih oblika u govoru paraćinskog Pomoravlja.} Belgrade},
	author = {Miloradović, Sofija},
	publisher = {Srpska akademija nauka i umetnosti},
	sortname = {Miloradovic, Sofija},
	year = {2003}
}


@book{Mirčev1958,
	address = {Sofia},
	author = {Mirčev, Kiril},
	publisher = {Nauka i izkustvo},
	sortname = {Mircev, Kiril},
	title = {\textit{Istoričeska gramatika na bâlgarskija ezik}},
	year = {1958}
}


@incollection{Montaut2018,
	address = {Berlin: Language Science Press. DOI},
	author = {Montaut, Annie},
	booktitle = {\textit{Diachrony of differential argument marking. (Studies in Diversity Linguistics 19)}},
	editor = {Ilja A. Seržant and Alena Witzlack-Makarevich},
	note = {5281/zenodo.1219168.},
	pages = {281--313},
	publisher = {10},
	title = {The rise of differential object marking in {Hindi} and related languages},
	year = {2018}
}


\textstyleauthor{Moravcsik, Edith A.}~\textstylepubyear{1978}.~\textstylechaptertitle{On the case marking of objects}. In~\textstyleeditor{J. H. Greenberg, C. A. Ferguson \& E. A. Moravcsik}( eds.).~\textstylebooktitle{\textit{Universals of human language: IV: Syntax,} }~\textstylepagefirst{249}{}-\textstylepagelast{289}. Stanford: Stanford University Press.

@incollection{Norde2002,
	address = {Amsterdam},
	author = {Norde, Muriel},
	booktitle = {\textit{Grammatical Relations in Change}},
	editor = {Jan Terje Faarlund},
	pages = {241--272},
	publisher = {John Benjamins},
	title = {The loss of lexical case in {Swedish}},
	year = {2002}
}


@book{Peco1985,
	address = {Belgrade},
	author = {Peco, Asim},
	publisher = {Naučna knjiga},
	title = {\textit{Pregled srpskohrvatskih dijalekata}},
	year = {1985}
}


@book{PiperPiper2013,
	address = {\textit{Normativna gramatika srpskog jezika.} Novi Sad},
	author = {Piper, Predrag and Ivan Klajn},
	publisher = {Matica Srpska},
	year = {2013}
}


@book{Popović1955,
	address = {\textit{Istorija srpskohrvatskog jezika.} Novi Sad},
	author = {Popović, Ivan},
	publisher = {Matica Srpska},
	sortname = {Popovic, Ivan},
	year = {1955}
}


@article{Rooth1992,
	author = {Rooth, Mats},
	journal = {\textit{Natural Language Semantics}},
	number = {1},
	pages = {75--116},
	title = {A Theory of Focus Interpretation},
	volume = {1},
	year = {1992}
}


@book{Rusek1964,
	address = {Wrocław},
	author = {Rusek, Jerzy},
	note = {Ossolińskich.},
	publisher = {Zakład Narodowy im},
	title = {\textit{Deklinacja i użycie przypadków w triodzie Chłudowa: {{S}}tudium nad rozwojem analityzmu w języku bułgarskim}},
	year = {1964}
}


@incollection{Schenker1993,
	address = {London},
	author = {Schenker, Alexander},
	booktitle = {\textit{The {Slavonic} Languages}},
	editor = {B. Comrie and G. G. Corbett},
	pages = {60--121},
	publisher = {Routledge},
	title = {Proto-{Slavic}},
	year = {1993}
}


@book{SeržantSeržant2018,
	address = {Berlin},
	booktitle = {\textit{Diachrony of differential argument marking}},
	editor = {Seržant, I., and Witzlack-Makarevich, A.},
	publisher = {Language Science Press},
	sortname = {Serzant, I., and Witzlack-Makarevich, A.},
	title = {\textit{Diachrony of differential argument marking}},
	year = {2018}
}


@incollection{Silverstein1976,
	address = {Canberra},
	author = {Silverstein, Michael},
	booktitle = {\textit{Grammatical categories in {Australian} languages}},
	editor = {R. M. W. Dixon},
	pages = {112--171},
	publisher = {Australian Institute of Aboriginal Studies},
	title = {Hierarchy of features and ergativity},
	year = {1976}
}


@book{Simić1972,
	address = {\textit{Srpski dijalektološki zbornik \REF{ex:key:19}.} Belgrade},
	author = {Simić, Radoje},
	publisher = {Srpska akademija nauka i umetnosti},
	sortname = {Simic, Radoje},
	title = {\textit{Levački govor}},
	year = {1972}
}


@book{Skautrup1944,
	address = {København},
	author = {Skautrup, Peter},
	publisher = {Gyldendal},
	title = {\textit{Det Danske Sprogs Historie I: {{F}}ra Guldhornene til Jyske Lov}},
	year = {1944}
}


@book{Slavova2017,
	address = {Sofia},
	author = {Slavova, Tatjana},
	note = {Kliment Ohridski” universitetsko izdatelstvo.},
	publisher = {“Sv},
	title = {\textit{Starobâlgarski ezik}},
	year = {2017}
}


\begin{styleBibliography}
Smith, Kenny, \& ~Jennifer Culbertson. 2020.~Communicative pressures shape language during communication (not learning): Evidence from case marking in artificial languages. https://doi.org/10.31234/osf.io/5nwhq.
\end{styleBibliography}

\begin{styleBibliography}
@article{Sobolev1991,
	author = {Sobolev, A.},
	journal = {\textit{Zbornik Matice srpske za filologiju i lingvistiku} XXXIV},
	pages = {93--139},
	title = {Kategoriia padeža na periferii balkanoslavjanskogo areala},
	volume = {/1},
	year = {1991}
}

\end{styleBibliography}

@book{Stanković2022,
	address = {\textit{Južnoslovenski filolog} 78 \REF{ex:key},
	author = {Stanković, Dragana},
	note = {49-76.},
	publisher = {1}},
	sortname = {Stankovic, Dragana},
	title = {Upotreba bespredloškog genitiva u govoru mlađih vranjanca},
	year = {2022}
}


@book{Stojkov1975,
	address = {Volume 3.  \textit{Jugozapadna Bâlgaria}. Sofia},
	booktitle = {\textit{Bâlgarski dialekten atlas}},
	editor = {Stojkov, Stojko},
	publisher = {Bâlgarska akademija na naukite},
	title = {\textit{Bâlgarski dialekten atlas}},
	year = {1975}
}


@book{Stojkov1981,
	address = {Volume. 4.  \textit{Severozapadna Bâlgaria}. Sofia},
	booktitle = {\textit{Bâlgarski dialekten atlas}},
	editor = {Stojkov, Stojko},
	publisher = {Bâlgarska akademija na naukite},
	title = {\textit{Bâlgarski dialekten atlas}},
	year = {1981}
}


@book{Todorov2002,
	address = {\textit{Starobâlgarskoto anaforično mestoimenie i negovite novobâlgarski zastâpnici.} Sofija},
	author = {Todorov, Todor},
	note = {Marin Drinov".},
	publisher = {"Prof},
	year = {2002}
}


@book{Tsonev1984,
	address = {Vol. 1. Sofija},
	author = {Tsonev, Benjo},
	publisher = {Nauka i izkustvo (2\textsuperscript{nd} edition)},
	title = {\textit{Istorija na bâlgarskija ezik}},
	year = {1984}
}


@misc{Valle2011,
	author = {Valle, Daniel},
	note = {\textit{Memorias del V Congreso de Idiomas Indígenas de Latinoamérica, 6-8 de octubre de 2011, Universidad de Texas en Austin}. \url{https://ailla.utexas.org/sites/default/files/documents/Valle_CILLA_V.pdf}},
	title = {Differential subject marking triggered by information structure},
	year = {2011}
}


\textstyleauthor{Witzlack{}-Makarevich, Alena}, \& Ilja A.~\textstyleauthor{Seržant}. \textstylepubyear{2018}.~\textstylechaptertitle{Differential argument marking: Patterns of variation}. In~\textstyleeditor{I. A. Seržant}~\& ~\textstyleeditor{A. Witzlack{}-Makarevich}~(eds.).~\textstylebooktitle{\textit{Diachronic typology of differential argument marking}}~. \textstylepagefirst{1}{}-\textstylepagelast{40}. Berlin: Language Science Press.

@article{Zimmermann2008,
	author = {Zimmermann, Malte},
	journal = {\textit{Acta Linguistica Hungarica}},
	pages = {347--360},
	title = {Contrastive focus and emphasis},
	volume = {3-4},
	year = {2008}
}


\end{verbatim}
\sloppy\printbibliography[heading=subbibliography,notkeyword=this]
\end{document}